\section{纤维化拓扑斯的投射极限}\etiquette[8]{VI.sec8}

\subsection{一般性质}\etiquette[8.1]{VI.subsec8.1}

\begin{enonce*}{定义 8.1.1}\etiquette[8.1.1]{VI.def8.1.1}
设 $𝒰$ 是一个宇宙, $p:F \to I$ 是一个纤维化 $𝒰$-拓扑斯.
由 $𝒰$-拓扑斯 $C$ 及 $I$ 上的纤维化 $𝒰$-拓扑斯的笛卡儿态射
$m : C \times I \to F$ 构成的对 $(C, m)$,
称为纤维化拓扑斯 $F$ 的投射极限,
是指对每个 $𝒰$-拓扑斯 $D$,
将基变换函子
$$
\Hhomtop(D, C) \to \Hhomtop_{I}(D\times I, C \times I)
$$
与同 $m$ 复合的函子 (\hyperref[VI.subsubsec7.1.8]{7.1.8})
复合而得到的函子
\begin{equation*}
\Hhomtop(D, C) \to \Hhomtop_{\cart/I}(D\times I, F)\tag{8.1.1.1}
\end{equation*}\etiquette[8.1.1.1]{VI.eq8.1.1.1}
为范畴等价.
\emph{(\cf \hyperref[VI.par8.1.3.3]{8.1.3.3} 中更“几何”的表述)}.
\end{enonce*}

\setcounter{subsubsection}{1}
\subsubsection{}\etiquette[8.1.2]{VI.subsubsec8.1.2}

沿用 \hyperref[VI.def8.1.1]{8.1.1} 的记号.
设 $D$ 是一个 $𝒰$-拓扑斯,
$m_1 : D \times I \to F$ 是纤维化 $𝒰$-拓扑斯的态射.
展开定义 \hyperref[VI.def8.1.1]{8.1.1},
立即可见存在拓扑斯态射 $n : D \to C$
及 $I$ 上的纤维化拓扑斯态射的同构
$\gamma : m_{1} \xrightarrow{\sim} m \circ (n \times \id_I)$,
使图交换:
\[\vcenter{
\xymatrix@C=.75cm@R=.5cm{
D \times I  \ar[rrrr]^{m_1} \ar[dddd]_{n\times \id_I} &&{\phantom{.}}\ar@/_.5pc/@{-->}[dr]_{\gamma}  &&F\\
                         & &{\phantom{a}}& &\\
                         & && & \\
                         & && & \\
C \times I \ar[rrrruuuu]_{m}&& &&{\phantom{a}}\\
}}\leqno{(8.1.2.1)}
\]\etiquette[8.1.2.1]{VI.eq8.1.2.1}
此外, 若 $(n', \gamma')$ 为另一个具有上述性质的对,
则存在\emph{唯一}同构 $\eta : n \xrightarrow{\sim} n'$,
使得
\begin{equation*}
\gamma \circ m (\eta \times \id) = \gamma'{\quoi}.\tag{8.1.2.2}
\end{equation*}\etiquette[8.1.2.2]{VI.eq8.1.2.2}
由此\pageoriginale{}可知,
若 $(D, m_1)$ 为 $F$ 的投射极限,
则 $n$ 为 $𝒰$-拓扑斯的等价
(IV \hyperref[IV.subsec3.4]{3.4});
使图 \eqhyperref[VI.eq8.1.2.1]{8.1.2.1} 交换的同构 $\gamma$ 的附加数据,
在典范同构意义下确定对 $(n, \gamma)$.
当 $I$ 余滤过时,
投射极限的存在性将在 \hyperref[VI.thm8.2.3]{8.2.3} 中证明.

\subsubsection{}\etiquette[8.1.3]{VI.subsubsec8.1.3}
以
$$
\limtop{I} F \mbox{ 或 } \limtop{I}
F_i \mbox{ 或又 } \underleftarrow{F}
$$
表示纤维化拓扑斯 $F$ 的一个投射极限,
以 $\mu : (\limtop{I} F) \times I \to F$ 表示典范态射,
并对 $I$ 的每个对象 $i$, 以
$$
\mu_i : \underleftarrow{F} \simeq \underleftarrow{F} \times {i} \to F_{i}
$$
表示 $\mu$ 在 $i$ 上的纤维上诱导的拓扑斯态射.

选定 $F$ 的一个双分裂（bisplitting）
(\hyperref[VI.subsubsec7.1.4]{7.1.4}),
\ie 对每个 $f : i \to j$ 选定拓扑斯态射
$f\bigdot = (f^{\ast}, f_{\ast}) : F_i \to F_j$,
使 $f^{\ast}$ 为纤维化结构的逆像函子.
于是对 $I$ 的每对可复合态射
$i \xrightarrow{f} j \xrightarrow{g} k$,
有拓扑斯态射的同构 (\hyperref[VI.subsubsec7.1.3]{7.1.3})
$$
c_{f, g} : g.f. \simeq (gf).\quad {\quoi}.
$$
由于 $\mu$ 是纤维化拓扑斯的笛卡儿态射,
对 $I$ 的每个态射 $f : i \to j$,
有拓扑斯态射的同构 (\hyperref[VI.subsubsec7.1.6]{7.1.6})
\begin{equation*}
b_f  : f. \mu_i \xrightarrow{\sim} \mu_j{\quoi},\tag{8.1.3.1}
\end{equation*}\etiquette[8.1.3.1.]{VI.eq8.1.3.1}
而各 $b_f$ 构成的族满足与
\eqhyperref[VI.eq7.1.6.2]{7.1.6.2} 类似的关系.

\setcounter{paragraph}{1}
\paragraph{}\etiquette[8.1.3.2]{VI.par8.1.3.2}

设 $D$ 是一个 $𝒰$-拓扑斯,
$m : D \times I \to F$ 是纤维化拓扑斯的态射,
$m_i : D = D \times \{i\} \to F_i$
为 $m$ 在各纤维上诱导的态射;
并对每个 $f : i \to j$, 设
$$
b'_f : f . m_i \xrightarrow{\sim} m_j{\quoi},
$$
为转移同构.
由 \hyperref[VI.subsubsec8.1.2]{8.1.2} 可知,
存在拓扑斯态射 $n : D \to \underleftarrow{F}$\pageoriginale{}
及一族同构
$$
\gamma_i : m_i \xrightarrow{\sim} \mu_i n \quad i \in \ob I {\quoi},
$$
使对每个 $f : i \to j$,
$$
b_f \circ f. (\gamma_i) = \gamma_j \circ b'_f{\quoi}.
$$
此外, 态射 $n$ 与同构族 $\gamma_i$,
$i \in \ob (I)$, 在唯一同构意义下确定
(\hyperref[VI.subsubsec8.1.2]{8.1.2}).

反过来, 假定给定拓扑斯 $\underleftarrow{F}$,
态射 $\mu_i : \underleftarrow{F} \to F_i$,
以及如 \eqhyperref[VI.eq8.1.3.1]{8.1.3.1.} 的同构 $b_f$,
它们满足通常的相容条件;
并假定所有这些数据具有上述泛性质.
则存在唯一的纤维化拓扑斯态射
$\mu : \underleftarrow{F} \times I \to F$,
给出各 $\mu_{i}$ 和 $b_f$,
而 $(\underleftarrow{F}, \mu)$ 为纤维化拓扑斯 $F$ 的投射极限.

\paragraph{}\etiquette[8.1.3.3]{VI.par8.1.3.3}

设 $𝒱$ 为一个满足 $\sheaf{\Uu} \in 𝒱$ 的宇宙.
上述讨论因而容许将纤维化拓扑斯的投射极限,
解释为取值于属于 $𝒱$ 的纤维化 $𝒰$-拓扑斯的 2-范畴的伪函子
在 2-范畴意义下的“投射极限”;
这些伪函子通过选取态射 $f.$ 从纤维化结构得到.
不过读者须注意,
不要将这种 2-范畴意义下的“投射极限”
(在同构意义下交换)
与投射极限的严格 (或朴素) 概念
(图严格交换) 混淆.
即使严格概念有意义
(当伪函子是真正的函子时),
广义“投射极限”与严格投射极限一般也不重合.

\subsubsection{}\etiquette[8.1.4]{VI.subsubsec8.1.4}

设 $p : F \to I$ 和 $q : G \to I$ 为两个纤维化 $𝒰$-拓扑斯,
$m : F \to G$ 为纤维化拓扑斯的态射.
假定纤维化拓扑斯 $F$ 与 $G$ 具有投射极限
(\hyperref[VI.def8.1.1]{8.1.1}).
由定义 \hyperref[VI.def8.1.1]{8.1.1} 立即可知,
存在拓扑斯态射
\begin{equation*}
\underleftarrow{m} : \underleftarrow{F} \to
\underleftarrow{G}{\quoi},\tag{8.1.4.1}
\end{equation*}\etiquette[8.1.4.1]{VI.eq8.1.4.1}
及纤维化拓扑斯态射的同构
$\gamma : m \circ \mu \xrightarrow{\sim}
\mu \circ \underleftarrow{m} \times \id_{I}$,
使图\pageoriginale{}交换:
\[\vcenter{
  \xymatrix@C=.6cm@R=.45cm{
\underleftarrow{F}\times I \ar[rrr]^{\mu}
\ar[ddd]_{\underleftarrow{m}\times \id_I}   && & F\ar[ddd]^{m}\\
                          &            & &\\
                          &            & &\ar@/_.55pc/@{-->}[ld]_{\gamma} \\
\underleftarrow{G}\times I \ar[rrr]^{\mu} &       & &  G&.
}}\leqno{(8.1.4.2)}
\]\etiquette[8.1.4.2]{VI.eq8.1.4.2}
此外, 若 $(\underleftarrow{m}', \gamma')$ 为另一个具有上述性质的对,
则存在\emph{唯一}拓扑斯态射的同构
$\eta: \underleftarrow{m} \xrightarrow{\sim} \underleftarrow{m}'$,
使得
$$
\gamma \circ \mu (\eta \times \id) = \gamma'{\quoi}.
$$

\subsubsection{}\etiquette[8.1.5]{VI.subsubsec8.1.5}

设 $p : F \to I$ 是属于 $𝒱$ 的一个纤维化 $𝒰$-拓扑斯,
$\phi : I'\to I$ 是属于 $𝒱$ 的函子,
$p_{\phi} : F_{\phi} \to I'$
为由 $(F, p)$ 经基变换 $\phi$ 得到的纤维化 $𝒰$-拓扑斯
(\hyperref[VI.subsubsec7.1.9]{7.1.9}).
设 $(\underleftarrow{F}, \mu)$ 为 $F$ 的一个投射极限.
由基变换, 得到 $I'$ 上的纤维化拓扑斯态射:
$$
\mu' : \underleftarrow{F} \times I' \to F_{\phi}{\quoi}.
$$
设 $(\underleftarrow{F}_{\phi}, \mu_{\phi})$
为 $F_{\phi}$ 的一个投射极限.
由投射极限的泛性质
(\hyperref[VI.subsubsec8.1.2]{8.1.2}),
于是有拓扑斯态射
$$
m_{\phi} : \underleftarrow{F} \to \underleftarrow{F}_{\phi}{\quoi},
$$
及 $I'$ 上的纤维化拓扑斯态射的同构:
$$
\gamma : \mu'\xrightarrow{\sim} \mu_{\phi} \circ (m_{\phi} \times \id_{I'}){\quoi}.
$$
对 $(m_{\phi}, \gamma)$ 在唯一同构意义下确定.

\subsubsection{}\etiquette[8.1.6]{VI.subsubsec8.1.6}

设 $p : F \to I$ 为小范畴 $I$ 上的一个纤维化 $𝒰$-景,
$F^{\sim/I} \to I$ 为相伴的纤维化 $𝒰$-拓扑斯
(\hyperref[VI.subsubsec7.2.6]{7.2.6}),
$D$ 为 $𝒰$-拓扑斯.
首先, 有函子
\begin{equation*}
\Hhomtop_{\cart / I} (D \times I, F^{\sim/I}) \to
\Mmorsite_{\cart/I}(D \times I, F){\quoi},\tag{8.1.6.1}
\end{equation*}\etiquette[8.1.6.1]{VI.eq8.1.6.1}
给每个 $m$ 对应 $m^{\ast}$ 与 $\epsilon_{F /I}$ 的复合
(\hyperref[VI.prop7.2.7]{7.2.7}).
此外, 由范畴 $\Mmorsite_{\cart /I} (D \times I, F)$ 的定义
\eqhyperref[VI.eq7.2.2.2]{7.2.2.2}, 有函子
\begin{equation*}
\Mmorsite_{\cart/I}(D \times I, F) \to \Hhom_{\cart / I} (F, D \times
I)^{\circ}{\quoi}.\tag{8.1.6.2}
\end{equation*}\etiquette[8.1.6.2]{VI.eq8.1.6.2}
最后\pageoriginale{}, 由归纳极限的定义
(\hyperref[VI.def6.3]{6.3}), 有函子
\begin{equation*}
\Hhom_{\cart / I} (F, D \times I)^{\circ} \to
\Hhom(\varinjlim_{I^{\circ}}F, D)^{\circ}{\quoi}.\tag{8.1.6.3}
\end{equation*}\etiquette[8.1.6.3]{VI.eq8.1.6.3}
以
\begin{equation*}
\Theta : \Hhomtop_{\cart / I}(D \times I, F^{\sim / I}) \to
\Hom(\varinjlim_{I^{\circ}}F, D)^{\circ}\tag{8.1.6.4}
\end{equation*}\etiquette[8.1.6.4]{VI.eq8.1.6.4}
表示最后这三个函子的复合,
并以
\begin{equation*}
\Pi : F \to \varinjlim_{I^{\circ}}F{\quoi},\tag{8.1.6.5}
\end{equation*}\etiquette[8.1.6.5]{VI.eq8.1.6.5}
表示典范函子.

\begin{enonce*}{命题 8.1.7}\etiquette[8.1.7]{VI.prop8.1.7}
函子 $\Theta$ 全忠实.
其本质像为那些使函子
$$
(m \circ \Pi, p) : F \to D \times I
$$
是从总景 $D \times I$ 到总景 $F$ 的态射的函子
$m : \varinjlim_{I^{\circ}}F \to D$ 所成的集合.
\end{enonce*}

由 \hyperref[VI.prop7.2.7]{7.2.7},
函子 \eqhyperref[VI.eq8.1.6.1]{8.1.6.1} 为范畴等价;
由 \hyperref[VI.subsubsec7.2.2]{7.2.2},
函子 \hyperref[VI.eq8.1.6.2]{(8.1.6.2)} 全忠实;
由 \hyperref[VI.def6.3]{6.3},
函子 \eqhyperref[VI.eq8.1.6.3]{8.1.6.3} 为范畴等价.
因此 $\Theta$ 全忠实.
函子 $m : \varinjlim_{I^{\circ}}F \to D$
同构于函子 $(m \circ \Pi, p)$ 在
$\Hhom_{\cart / I} (F, D \times I)^{\circ} \to
\Hhom(\varinjlim_{I^{\circ}}F, D)^{\circ}$ 下的像;
这样的函子属于 \eqhyperref[VI.eq8.1.6.2]{8.1.6.2} 的本质像,
当且仅当它为从总景 $D \times I$ 到总景 $F$ 的态射
(\hyperref[VI.thm7.4.13.1]{7.4.13.1});
从而得到命题.

% END INCLUDED SGA4_English_current_buildable_source_20260906/source/06/drafts/SGA4_VI_LINES_5252_5495_English.texfrag


% BEGIN INCLUDED SGA4_English_current_buildable_source_20260906/source/06/drafts/SGA4_VI_LINES_5496_5856_English.texfrag
\subsection{指标范畴余滤过时投射极限的构造}\etiquette[8.2]{VI.subsec8.2}

\begin{enonce*}{引理 8.2.1}\etiquette[8.2.1]{VI.lem8.2.1}
设 $I$ 与 $I'$ 是两个余滤过范畴
\emph{(I \hyperref[I.sec8.1]{8.1})},
$\phi:I'\to I$ 为一个使
$\phi^{\circ}:{I'}^{\circ}\to I^{\circ}$ 共尾的函子
\emph{(I \hyperref[I.sec8.1]{8.1})}.
再设 $p:F\to I$ 是纤维化 $𝒰$-拓扑斯,
$D$ 是 $𝒰$-拓扑斯,
$m:D\times I\to F$ 为纤维化拓扑斯的笛卡儿态射,
$m':D\times I'\to F'$ 为由 $m$ 经基变换 $\phi$ 得到的纤维化拓扑斯态射
(\hyperref[VI.subsubsec7.1.9]{7.1.9}).
对 $(D,m)$ 是 $F$ 的投射极限,
当且仅当 $(D,m')$ 是 $F'$ 的投射极限.
\end{enonce*}

证明\pageoriginale{}只是相当冗长的例行验证.
留给耐心的读者,
可采用伪函子投射极限的观点
(\hyperref[VI.par8.1.3.3]{8.1.3.3}).

\begin{enonce*}{引理 8.2.2}\etiquette[8.2.2]{VI.lem8.2.2}
设 $p:F\to I$ 是小余滤过范畴上的纤维化 $𝒰$-景,
$\underrightarrow{F}$ 是纤维化范畴 $F$ 的归纳极限
\hyperref[VI.def6.3]{(6.3.9)},
$\pi:F\to\underrightarrow{F}$ 是典范函子.
给 $\underrightarrow{F}$ 配备使 $\pi$ 连续的最粗拓扑
\emph{(III 1)}.
设 $D$ 为 $𝒰$-景,
$n:\underrightarrow{F}\to D$ 为函子.
下列性质等价:
\begin{enumerate}
\renewcommand{\theenumi}{\roman{enumi}}
\renewcommand{\labelenumi}{\theenumi)}
\item $n$ 为景态射 $D\to\underrightarrow{F}$.

\item $n\circ\pi$ 为景态射
(其中 $F$ 配备总拓扑).

\item $(n\circ\pi,p):F\to D\times I$
为从总景 $D\times I$ 到总景 $F$ 的态射.

\item 对 $I$ 的每个对象 $i$,
复合函子 $F_i\xrightarrow{\alpha_{i!}}F\xrightarrow{n\circ\pi}D$
为景态射.
\end{enumerate}
\end{enonce*}

 i) 与 ii) 的等价性来自 \hyperref[VI.exe6.6]{6.6},
 iii) 与 iv) 的等价性来自 \hyperref[VI.thm7.4.13.1]{7.4.13}.
证明 iii) 与 ii) 的等价性.
为此只须证明第一投影函子 $\pr_1:D\times I\to D$ 为景态射.
以 $D^{\sim}$ 表示 $D$ 上的层拓扑斯.
由 \hyperref[VI.prop7.4.7]{7.4.7},
$(D\times I)^{\sim}$ 等价于拓扑斯
$\Hhom_{I^{\circ}}(I^{\circ},D^{\sim}\times I^{\circ})=
\Hhom(I^{\circ},D^{\sim})$.
此外, 立即可见 $\pr_1:D\times I\to D$
到层的自然扩张 (III \hyperref[III.prop1.2]{1.2})
是从 $\Hhom(I^{\circ},D^{\sim})$ 到 $D^{\sim}$ 的函子,
它给函子 $i\mapsto\sigma(i)$
对应对象 $\varinjlim_{I^{\circ}}\sigma(i)$.
由于 $I^{\circ}$ 滤过,
关于 $I^{\circ}$ 的归纳极限正合 (II 4),
故 $\pr_1$ 为景态射 (IV \hyperref[IV.subsec4.9]{4.9}).
还须证明 ii) 与 iv) 的等价性.
先处理 $D$ 是 $𝒰$-拓扑斯,
而 $F$ 是 $I$ 上的纤维化 $𝒰$-拓扑斯的情形.

设 $i$ 是 $I$ 的对象.
仿照 \hyperref[VI.par7.4.13.6]{7.4.13.6} 论证,
可见函子 $\alpha_{i!}:F_i\to F$ 分解为态射
$\bar{\alpha}_{i!}:F_i\to F/\alpha_{i!}(e_i)$
后接局部化函子 $F/\alpha_{i!}(e_i)\to F$,
其中 $e_i$ 为 $F_i$ 的终对象.
因此只须证明复合函子
$F/\alpha_{i!}(e_i)\xrightarrow{n\circ\pi}D$ 为景态射.
但存在在同构意义下交换的图:
\[
\xymatrix@C=2.5cm@R=1.25cm{
F/\alpha_{i!}(e_i) \ar[d]\ar[r]^{n\circ
  \pi/\alpha_{i!}(e_i)}& D/n\circ \pi \alpha_{i!}(e_i) \ar[d]\\
F \ar[r]^{n\circ \pi}& D
}
\]
由于\pageoriginale{}作为景态射的性质是局部的
(IV \hyperref[IV.subsec5.10]{5.10}),
只须证明 $n\circ\pi\circ\alpha_{i!}(e_i)$ 为 $D$ 的终对象.
对 $I$ 的每个态射 $f:j\to k$,
在 $f$ 上存在唯一态射
$e_f:\alpha_{j!}(e_j)\to\alpha_{k!}(e_k)$,
且此态射为笛卡儿的.
因此 $j\mapsto\alpha_{j!}(e_j)$
是 $I$ 上的 $F$ 的笛卡儿截面.
故典范态射
$n\circ\pi\circ\alpha_{i!}(e_i)\to
\varinjlim_I n\circ\pi\circ\alpha_{j!}(e_j)$ 为同构,
因为 $\pi$ 将笛卡儿态射映为同构.
此外, $\varinjlim_I\alpha_{j!}(e_j)$
(其中归纳极限在 $F^{\widehat{\phantom{,}}}$ 中取)
为 $F^{\widehat{\phantom{,}}}$ 的终对象.
由于 $n\circ\pi$ 为景态射,
它将 $F^{\sim}$ 的终对象映为 $D$ 的终对象.
因此 $\varinjlim_I n\circ\pi\circ\alpha_{j!}(e_j)$
为 $D$ 的终对象, 这就完成此情形的证明.

在一般情形, 函子 $(n\pi,p):F\to D\times I$
为笛卡儿且连续的 (\hyperref[VI.prop7.4.4]{7.4.4}).
此外 $n\circ\pi$ 为复合函子
$$
F \xrightarrow{(n \pi, p)} D \times I \xrightarrow{\pr_1} D{\quoi}.
$$
过渡到相伴纤维化 $𝒰$-拓扑斯,
得到在同构意义下交换的图:
\[
\xymatrix@C=2cm@R=1.75cm{
F       \ar@{-}[r]^{(n\pi,p)} \ar[d]_{\varepsilon_{F/I}}& D\times I
\ar[r]^{\pr_1}\ar[d]^{\varepsilon_D\times \id} & D \ar[d]^{\varepsilon_D}\\
F^{\sim/I}  \ar[r]^{(n\pi,p)^{\sim/I}} & D^{\sim}\times I \ar[r]^{\pr_1} & D^{\sim}&.
}
\]
由于 $(n\pi,p)^{\sim/I}$ 为笛卡儿的,
它形如 $(n'\pi',p')$,
其中 $p':F^{\sim/I}\to I$ 及
$\pi':F^{\sim/I}\to\varinjlim_{I^{\circ}}F^{\sim}_i$
是典范函子.
因此 $\pr_{1^{\circ}}(n\pi,p)^{\sim/I}=n'$.
由于函子 $\epsilon_{F/I}$ 与 $\epsilon_D$
在层拓扑斯上诱导等价 (\hyperref[VI.cor7.4.11]{7.4.11}),
函子 $n$ 与 $n'$ 在过渡到相伴层后诱导同构的函子
(III \hyperref[III.prop1.2]{1.2}),
故 $n'$ 为景态射.
由前面所证,\pageoriginale{}
函子 $(n\pi,p)^{\sim/I}$ 在各纤维上诱导拓扑斯态射.
故它为从总景 $F^{\sim/I}$ 到总景 $D^{\sim}\times I$ 的态射
\hyperref[VI.thm7.4.13.1]{(7.4.13)}.
由于函子 $\epsilon_{F/I}$ 与 $\epsilon_D\times\id$
在总拓扑斯上诱导等价 (\hyperref[VI.cor7.4.11]{7.4.11}),
函子 $(n\pi,p)$ 为纤维化景态射
\hyperref[VI.thm7.4.13.1]{(7.4.13)}, 这就证明了断言.

\begin{enonce*}{定理 8.2.3}\etiquette[8.2.3]{VI.thm8.2.3}
设 $p:F\to I$ 是一个本质小余滤过范畴
\emph{(I \hyperref[I.sec8.1]{8.1})} 上的纤维化 $𝒰$-景.
\begin{enumerate}
\renewcommand{\labelenumi}{\theenumi)}
\item 以 $\underrightarrow{F}$ 表示底范畴为纤维化范畴 $F\to I$ 的归纳极限
\emph{(\hyperref[VI.def6.3]{6.3})} 的景,
其拓扑为使典范函子 $\pi:F\to\underrightarrow{F}$ 连续的最粗拓扑.
则 $\underrightarrow{F}$ 是 $𝒰$-景,
且
$$
(\pi, p) : F \to \underrightarrow{F} \times I
$$
为从纤维化景 $\underrightarrow{F}\times I$ 到纤维化景 $F$ 的态射.

\item 设 $\mu:\underrightarrow{F}^{\sim}\times I\to F^{\sim/I}$
为由 $(\pi,p)$ 过渡到相伴纤维化 $𝒰$-拓扑斯
所得的纤维化 $𝒰$-拓扑斯态射.
对 $(\underrightarrow{F}^{\sim},\mu)$
是 $F^{\sim/I}$ 的投射极限.
\end{enumerate}
\end{enonce*}

\begin{enonce*}[remark]{注 8.2.4}\etiquette[8.2.4]{VI.rem8.2.4}
\begin{enumerate}
\renewcommand{\labelenumi}{\theenumi)}
\item 根据 \hyperref[VI.subsubsec8.1.3]{8.1.3} 引入的记号, 可置
$$
\limtop{I} F^{\sim} = \underrightarrow{F}^{\sim}.
$$

\item 由 \hyperref[VI.prop7.4.4]{7.4.4} 可知,
$\underrightarrow{F}$ 上的拓扑
也是使函子族 $\pi_j:F_j\to F\xrightarrow{\pi}\underrightarrow{F}$,
$j\in\ob I$, 连续的最粗拓扑.
\end{enumerate}
\end{enonce*}

\begin{enonce*}{定义 8.2.5}\etiquette[8.2.5]{VI.def8.2.5}
\hyperref[VI.thm8.2.3]{8.2.3} 引入的景 $\underrightarrow{F}$,
称为纤维化景 $F$ 的投射极限景.
\end{enonce*}

作为练习, 读者可仿照 \hyperref[VI.lem8.2.1]{8.2.1},
定义纤维化景的投射极限概念,
并验证景 $\underrightarrow{F}$ 确为景 $F$ 的投射极限.
因此, 定理的第二个断言可以改述为:
投射极限景上的层拓扑斯等价于相伴纤维化拓扑斯的投射极限.

\setcounter{subsubsection}{5}
\subsubsection{定理的证明: 归结为小指标范畴的情形}\etiquette[8.2.6]{VI.subsubsec8.2.6}

我们\pageoriginale{}只给提示.
设 $I'$ 为 $I$ 的满子范畴,
使 ${I'}^{\circ}$ 在 $I^{\circ}$ 中共尾.
以 $F'\to I'$ 表示由 $F$ 经基变换 $I'\to I$ 得到的纤维化范畴,
以 $\pi':F'\to\underrightarrow{F}'$ 表示典范函子.
先观察到范畴 $\underrightarrow{F}$ 与 $\underrightarrow{F}'$ 等价,
且经此等价从 $\underrightarrow{F}$ 的拓扑得到的
$\underrightarrow{F}'$ 上的拓扑,
是使函子 $\pi':F'\to\underrightarrow{F}'$ 连续的最粗拓扑.
定理的断言 1) 一旦在 $I$ 小的情形得证, 就随之成立.
一般情形的断言 2), 再由 $I$ 小的情形的断言 2)
及引理 \hyperref[VI.lem8.2.1]{8.2.1} 得到.

\subsubsection{定理证明的结束}\etiquette[8.2.7]{VI.subsubsec8.2.7}
第一个断言来自引理 \hyperref[VI.lem8.2.2]{8.2.2},
其中取 $D=\underrightarrow{F}$ 及 $n=\id$.
设 $D$ 是一个 $𝒰$-拓扑斯.
有在同构意义下交换的图:
{\def\objectstyle{\scriptstyle}
\def\labelstyle{\scriptstyle}
\[
\xymatrix@C=.8cm@R=1.15cm{
\Hhomtop(D,{(\underrightarrow{F})}^{\sim}) \ar[r]^{(1)}\ar[d]_{(3)}&
\Mmorsite(D,\underrightarrow{F}) \ar@{^{(}->}[r]^{(2)} \ar[d]_{(3)}&
\Hhom(\underrightarrow{F},D)^{\circ}\ar[dd]_{(5)}\\
\Hhomtop_{\Cart/I}(D\times I,(\underrightarrow{F})^{\sim}\times
I)  \ar[r]^{(1)} \ar[d]_{(4)}& \Mmorsite_{\Cart/I}(D\times I,
(\underrightarrow{F})\times I) \ar[d]&\\
\Hhomtop_{\Cart/I}(D\times I,F^{\sim/I})
\ar[r]^{(1)} & \Mmorsite_{\Cart/I}(D\times I,F)  \ar@{^{(}->}[r]^{(2)} &
\Hhom_{\Cart/I}(F,D\times I)^{\circ},
}
\]}

\noindent 其中类型 (1) 的函子为等价
(IV \hyperref[IV.prop4.9.4]{4.9.4} 及 \hyperref[VI.prop7.2.7]{7.2.7}),
类型 (2) 的函子全忠实
(IV \hyperref[IV.subsubsec4.9.1]{4.9.1} 及 \hyperref[IV.eq7.2.2]{7.2.2}),
类型 (3) 的函子为基变换函子,
类型 (4) 的函子为复合函子
(与 $\mu$, \resp $(\pi,p)$ 复合),
而函子 (5) 为由归纳极限的泛性质
(\hyperref[VI.def6.3]{6.3}) 得到的等价.
要证明函子
$\Hhomtop(D,(\underrightarrow{F})^{\sim})\xrightarrow{(4)\circ(3)}
\Hhomtop_{/I}(D\times I,F^{\sim/I})$ 为等价,
只须证明函子
$$
\Mmorsite(D, \underrightarrow{F})
\xrightarrow{(4) \circ (3)} \Mmorsite_{/I}(D \times I,F)
$$
为范畴\pageoriginale{}等价.
等价地, 只须证明两个全忠实函子
\begin{align*}
& \Mmorsite(D, \underrightarrow{F}) \xrightarrow{(5) \circ (2)} \Hhomcart_I(F, D \times I)^{\circ}\\
& \Mmorsite_{/I}(D \times I, F) \xrightarrow{(2)} \Hhomcart_I(F, D \times
I)^{\circ}
\end{align*}
具有相同本质像,
这立即来自引理 \hyperref[VI.lem8.2.2]{8.2.2}.

\subsubsection{}\etiquette[8.2.8]{VI.subsubsec8.2.8}

设 $p:F\to I$ 是范畴 $I$ 上的纤维化 $𝒰$-拓扑斯,
$(\limtop{I}F_i,\mu)$ 是 $F$ 的一个投射极限.
因此有 $I$ 上的纤维化 $𝒰$-拓扑斯的笛卡儿态射:
$$
\mu:\limtop{I} F_i \times I \to F{\quoi},
$$
从而, 通过直像函子过渡到 $I^{\circ}$ 上的纤维化范畴后
(\hyperref[VI.subsubsec7.1.3]{7.1.3}), 有笛卡儿函子
$$
\mu_{\ast}: \limtop{I} F_i \times I^{\circ} \to F'
$$
它是 $\mu$ 的直像函子 (\hyperref[VI.subsubsec7.1.7]{7.1.7}).
因此, 由纤维化范畴的投射极限的定义
(\hyperref[VI.subsec6.10]{6.10}), 有典范函子
\begin{equation*}
\Theta : \limtop{I} F_i  \to
\varprojlim_{I, f_{\ast}} F_i = \Hhom_{\cart/I^{\circ}} (I^{\circ},
F'){\quoi}.\tag{8.2.8.1}
\end{equation*}\etiquette[8.2.8.1]{VI.eq8.2.8.1}

\begin{enonce*}{定理 8.2.9}\etiquette[8.2.9]{VI.thm8.2.9}
沿用 \hyperref[VI.subsubsec8.2.8]{8.2.8} 的记号,
并假定 $I$ 为本质小余滤过范畴.
则函子 $\Theta$ 为范畴等价.
假定 $I$ 小. 置 $\limtop{I}F_i=(\underrightarrow{F})^{\sim}$
(由 \hyperref[VI.thm8.2.3]{8.2.3}
及 \hyperref[VI.rem8.2.4]{8.2.4}, 可如此),
并将范畴 $\Hhom_{I^{\circ}}(I^{\circ},F')$
解释为总景 $F$ 上的层范畴 (\hyperref[VI.prop7.4.7]{7.4.7}).
则复合函子
$$
(\underrightarrow{F})^{\sim} = \limtop{I} F_i
\xrightarrow{\Theta} \Hhom_{\cart/I^{\circ}} (I^{\circ}, F') \hookrightarrow \Hhom_{I^{\circ}}
(I^{\circ}, F') \approxeq F^{\sim}{\quoi},
$$
正是由函子
$$
\pi : F \to \underrightarrow{F}{\quoi}.
$$
定义的景态射的直像函子.
\end{enonce*}

\setcounter{subsubsection}{9}
\subsubsection{}\etiquette[8.2.10]{VI.subsubsec8.2.10}

在证明定理 \hyperref[VI.thm8.2.9]{8.2.9} 前,
先作一些\pageoriginale{}说明.
除了解决指标范畴余滤过时投射极限存在的问题
(在这一理论中, 此问题的解答颇为非平凡),
定理 \hyperref[VI.thm8.2.3]{8.2.3}
及 \hyperref[VI.thm8.2.9]{8.2.9}
还给出实际使用投射极限拓扑斯时所必需的另外两项信息.
首先, 定理 \hyperref[VI.thm8.2.3]{8.2.3}
将投射极限拓扑斯解释为某个景上的层拓扑斯,
该景的底范畴为各范畴 $F_i$ 在逆像函子系统下的归纳极限范畴.
在应用中, 可以具体解释此范畴及其拓扑.
其次, 定理 \hyperref[VI.thm8.2.9]{8.2.9}
断言投射极限上的层由它在各拓扑斯 $F_i$ 中的直像系统确定
(它给出范畴 $(F^{\sim/I})'$ 在 $I^{\circ}$ 上的一个笛卡儿截面).
反过来, 定理还断言:
若给定 $I$ 的每个对象 $i$ 上的 $F_i$ 的层 $X_i$,
并对 $I$ 的每个态射 $f:i\to j$ 给定同构
$X_j\simeq f_{\ast}X_i$,
这些同构满足相容条件,
则投射极限上本质上存在唯一的层,
其直像给出各 $X_i$ 的系统.

\subsubsection{}\etiquette[8.2.11]{VI.subsubsec8.2.11}

\hyperref[VI.thm8.2.3]{8.2.3} 的断言 2)
及 \hyperref[VI.thm8.2.9]{8.2.9} 的第一个断言,
可概括为引人注目的公式
\begin{equation*}
\varprojlim_{I, f_{\ast}}F_i^{\sim} \approxeq
\limtop{I, f.} F_i^{\sim} \approxeq
(\varinjlim_{I^{\circ}, f_{\ast}} F_i)^{\sim}{\quoi}. \tag{8.2.11.1}
\end{equation*}\etiquette[8.2.11.1]{VI.eq8.2.11.1}

\subsubsection{}\etiquette[8.2.12]{VI.subsubsec8.2.12}
\emph{定理 \hyperref[VI.thm8.2.9]{8.2.9} 的证明}.
要证明第一个断言,
利用 \hyperref[VI.lem8.2.1]{8.2.1}
以及关于纤维化范畴投射极限的类似引理,
归结为范畴 $I$ 小的情形.
细节留给读者.
以下假定 $I$ 是小范畴.
已知函子
$$
\pi : F \to \varinjlim_{I^{\circ}} F_i
$$
为景态射 (\hyperref[VI.lem8.2.2]{8.2.2}).
它在拓扑斯上定义的\pageoriginale{}直像函子
$$
\pi_{\ast} : (\varinjlim_{I^{\circ}} F_i)^{\sim} \to F^{\sim}
$$
于是全忠实,
其本质像是 $F$ 上将 $F$ 的笛卡儿态射映为同构的层所成的集合
(\hyperref[VI.exe6.6]{6.6}).
在等价 $F^{\sim}\simeq\Hhom_{I^{\circ}}(I^{\circ},F')$ 下,
这些层对应于 $F'$ 在 $I^{\circ}$ 上的笛卡儿截面
(\hyperref[VI.prop7.4.7]{7.4.7}).
此外, 对每个 $j\in\ob(I)$,
以 $\alpha_{j!}:F_j\to F$ 表示包含函子,
它给出三个函子
$\xymatrix{(\alpha_{j!},\alpha^{\ast}_j,\alpha_{j\ast}):F_j
\ar@<4pt>[r]\ar@<-4pt>[r]&F\ar[l]}$
(\hyperref[VI.rem7.4.3]{7.4.3}).
由等价 $F^{\sim}\simeq\Hhom_{I^{\circ}}(I^{\circ},F')$ 的定义
(\hyperref[VI.prop7.4.7]{7.4.7}) 可知,
$\varinjlim_{I^{\circ}}F_i$ 上的层 $X$
在函子
$(\varinjlim_{I^{\circ}}F_i)^{\sim}\xrightarrow{\pi_{\ast}}F^{\sim}
\simeq\Hhom_{I^{\circ}}(I^{\circ},F')$ 下
映为笛卡儿截面 $j\mapsto\alpha^{\ast}_j\pi_{\ast}(X)$.
但 \hyperref[VI.thm8.2.3]{8.2.3} 对函子
$\mu:(\varinjlim_{I^{\circ}}F_i)^{\sim}\times I\to F$ 的描述表明,
对 $I$ 的每个对象 $j$,
由 $\mu$ 过渡到 $j$ 处纤维得到的态射
$\mu_j:(\varinjlim_{I^{\circ}}F_i)^{\sim}\to F_j$
的直像函子
$\mu_{j\ast}:(\varinjlim_{I^{\circ}}F_i)^{\sim}\to F_j$
正是函子 $\alpha^{\ast}_j\pi_{\ast}$.
因此, 函子
$(\varinjlim_{I^{\circ}}F_i)^{\sim}\xrightarrow{\pi_{\ast}}F^{\sim}
\simeq\Hhom_{I^{\circ}}(I^{\circ},F')$
正是函子
$(\varinjlim_{I^{\circ}}F_i)^{\sim}\xrightarrow{\Theta}
\Hhom_{\cart/I^{\circ}}(I^{\circ},F')\hookrightarrow
\Hhom_{I^{\circ}}(I^{\circ},F')$,
这就证明了定理.


% END INCLUDED SGA4_English_current_buildable_source_20260906/source/06/drafts/SGA4_VI_LINES_5496_5856_English.texfrag


% BEGIN INCLUDED SGA4_English_current_buildable_source_20260906/source/06/drafts/SGA4_VI_LINES_5857_6042_English.texfrag
\subsection{投射极限景的拓扑: 凝聚拓扑斯的情形}\etiquette[8.3]{VI.subsec8.3}

\subsubsection{}\etiquette[8.3.1]{VI.subsubsec8.3.1}

本节中, $p:F\to I$ 是本质小余滤过范畴上的纤维化 $\Uu$-景,
其逆像函子 $f^{\ast}:F_i\to F_j$ 与纤维积交换.
此外, 对每个纤维范畴 $F_i$ 的每个对象 $X$,
给定 $F_i$ 的拓扑的覆盖族的一个集合 $\Cov(X)$,
具有 II \hyperref[II.def1.3]{1.3} 的性质 (PTO) 与 (PT1),
并生成 $F_i$ 的拓扑.
最后, 假定对每个 $f:i\to j$、每个 $Y\in\ob F_j$
及 $\Cov(Y)$ 中的每个族 $(Y_{\alpha}\to Y)_{\alpha\in A}$,
族 $(f^{\ast}(Y_{\alpha})\to f^{\ast}(Y))_{\alpha\in A}$
属于 $\Cov(f^{\ast}(Y))$.
例如取 $\Cov(X)=$ \{$F_i$ 中 $X$ 的所有覆盖族\},
这些关于 $\Cov(X)$ 的条件就满足.
本节\pageoriginale{}将在适当的有限性假设下
研究 $F$ 的投射极限景
(\hyperref[VI.thm8.3.13]{8.3.13}).

\subsubsection{}\etiquette[8.3.2]{VI.subsubsec8.3.2}

设 $\pi:F\to\underrightarrow{F}$ 为典范函子,
$Y$ 是 $\underrightarrow{F}$ 的一个对象.
以 $\Cov(Y)$ 表示所有具有如下形式的族
$(Y_{\alpha}\xrightarrow{n_{\alpha}}Y)_{\alpha\in A}$ 所成的集合:

存在 $i\in I$、$F_i$ 的一个对象 $X$、族
$(X_{\alpha}\xrightarrow{{n'}_\alpha}X)_{\alpha\in A}\in\Cov(X)$、
同构 $\phi:\pi(X)\simeq Y$
及同构族 $\phi_{\alpha}:\pi(X_{\alpha})\simeq Y_{\alpha}$,
$\alpha\in A$,
使对每个 $\alpha\in A$, 图
\[
\xymatrix@C=2.5cm@R=2cm{
\pi(X_{\alpha})
\ar@{-}[r]^{\displaystyle{\mathop{\sim}^{\phi_{\alpha}}}} \ar[d]^{\pi(n'_{\alpha})}& Y_{\alpha}\ar[d]^{n_{\alpha}} \\
\pi(X)  \ar@{-}[r]^{\displaystyle{\mathop{\sim}^{\phi}}} & Y
}
\]
交换.

\begin{enonce*}{命题 8.3.3}\etiquette[8.3.3]{VI.prop8.3.3}
\begin{enumerate}
\renewcommand{\labelenumi}{\theenumi)}
\item 对各 $Y\in\ob\underrightarrow{F}$,
属于 $\Cov(Y)$ 的各族生成投射极限景的拓扑
(\hyperref[VI.def8.2.5]{8.2.5}).

\item 族 $Y\mapsto\Cov(Y)$, $Y\in\ob\underrightarrow{F}$,
具有 \emph{II \hyperref[II.def1.3]{1.3}}
的性质 \emph{(PTO)} 与 \emph{(PT1)}.
\end{enumerate}
\end{enonce*}

将需要下述引理.

\begin{enonce*}{引理 8.3.4}\etiquette[8.3.4]{VI.lem8.3.4}
设 $i$ 为 $I$ 的对象,
$n:X'\to X$ 为 $F_i$ 的可作基变换的态射,
且对每个 $f:j\to i$, $f^{\ast}(n)$ 可作基变换.
则 $\pi(n)$ 可作基变换.
各函子
$$
\pi_j:F_j\hookrightarrow F\xrightarrow{\pi}\underrightarrow{F},
\quad j\in\ob I{\quoi},
$$
与纤维积交换.
\end{enonce*}

回忆, $\underrightarrow{F}$ 的每个对象 $Y$
都等于某个对象 $\pi(Z)$, $Z\in\ob F$,
而每个态射 $m:Y\to\pi(X)$ 都形如
$$
\pi(Z)\xrightarrow{\pi(s)^{-1}}\pi(Z')
\xrightarrow{\pi(m')}\pi(X){\quoi},
$$
其中 $s$ 为 $F$ 的笛卡儿态射.
因此, 要证明 $\pi(n)$ 可作基变换,\pageoriginale{}
只须证明每个图
\[
\xymatrix@C=2.5cm@R=1.25cm{
                           &\pi(X')\ar[d]^{\pi(n)}\\
\pi(Z')\ar[r]^{\pi(m')}&\pi(X)
}
\]
的纤维积可表示.
态射 $m':Z'\to X$ 分解为
$Z'\xrightarrow{m''}X''\xrightarrow{s'}X$,
其中 $p(m'')$ 为 $p(Z')=j$ 的恒等,
而 $s'$ 在 $p(s')=f$ 上为笛卡儿的.
置 $X'''=f^{\ast}(X')$ 及 $n'=f^{\ast}(n)$.
于是有交换图
\[
\xymatrix@C=2.5cm@R=1.5cm{
X'''\ar[r]^{s''}\ar[d]^{n'}&X'\ar[d]^{n}\\
X''\ar[r]^{s'}&X,
}
\]
从而在 $\underrightarrow{F}$ 中有交换图:
\[
\xymatrix@C=2.5cm@R=1.5cm{
        &\pi(X''')\ar[r]^{\sim}\ar[d]_{\pi(n')}&\pi(X')\ar[d]^{\pi(n)}\\
\pi(Z')\ar@{-}[r]^{\pi(m'')}&\pi(X'')\ar[r]^{\sim}&\pi(X)
}
\]
其中 $n'$ 为 $F_j$ 的可作基变换的态射.
因此, 要证明 $\pi(n)$ 可作基变换,
只须证明积 $\pi(Z')\times_{\pi(X'')}\pi(X''')$ 可表示;
所以只须证明 $\pi_j:F_j\hookrightarrow F\to\underrightarrow{F}$
与纤维积交换.

设
\[
\xymatrix@C=2.5cm@R=1.25cm{
Y'\ar[r]\ar[d]&X'\ar[d]\\
Y\ar[r]&X
}
\]
是 $F_j$ 中的笛卡儿图,
$W$ 为 $F$ 中位于 $k\in\ob I$ 上的对象.
有 (\hyperref[VI.prop6.5]{6.5}):
\def\labelstyle{\scriptstyle}
$$\displaystyle{\Hhom(\pi(W),\pi(Y'))\simeq\mathop{\varinjlim_{
\xymatrix@R=.3cm@C=.3cm@M=.05cm{
            &j\\
        l\ar[ru]^{f}\ar[rd]_{g}&\\
            &k
}}}\Hhom_{F_{\ell}}(g^{\ast}(W),f^{\ast}(Y'))}$$

由于\pageoriginale{} $f^{\ast}:F_j\to F_{\ell}$ 与纤维积交换, 有
$$
\Hom(g^{\ast}(W),f^{\ast}(Y'))\simeq
\Hom(g^{\ast}(W),f^{\ast}(Y))
\times_{\Hom(g^{\ast}(W),f^{\ast}(X))}
\Hom(g^{\ast}(W),f^{\ast}(X')).
$$
利用滤过极限与纤维积的交换 (I 2)
及公式 (\hyperref[VI.prop6.5]{6.5}), 得到同构
$$
\Hom(\pi(W),\pi(Y'))\simeq\Hom(\pi(W),\pi(Y))
\times_{\Hom(\pi(W),\pi(X))}\Hom(\pi(W),\pi(X')){\quoi},
$$
这表明图
\[
\xymatrix@C=2.5cm@R=1.5cm{
\pi(Y')\ar[r]\ar[d]&\pi(X')\ar[d]\\
\pi(Y)\ar[r]&\pi(X)
}
\]
在 $\underrightarrow{F}$ 中为笛卡儿图.

\setcounter{subsubsection}{4}
\subsubsection{}\etiquette[8.3.5]{VI.subsubsec8.3.5}
\emph{\hyperref[VI.prop8.3.3]{8.3.3} 的证明}.
先证明第二个断言.
由 \hyperref[VI.lem8.3.4]{8.3.4} 可知,
各 $\Cov(Y)$, $Y\in\underrightarrow{F}$,
中的族由可作基变换的态射构成, 从而得到性质 (PTO).
要证明这些族具有性质 (PT1) (基变换下稳定性),
经过与 \hyperref[VI.prop8.3.3]{8.3.3} 的证明中类似的一系列归结,
只须证明下列断言:

对每个 $i\in\ob I$、每个 $Y\in\ob F_i$、每个
$(Y_{\alpha}\to Y)_{\alpha\in A}\in\Cov(Y)$
及 $F_i$ 的每个态射 $u:Z\to Y$, 族
$$
(\pi(Z)\times_{\pi(Y)}\pi(Y_{\alpha})\to\pi(Z))_{\alpha\in A}
$$
属于 $\Cov(\pi(Z))$.

这个断言来自各
$\pi_i:F_i\hookrightarrow F\xrightarrow{\pi}\underrightarrow{F}$
与纤维积交换
(\hyperref[VI.lem8.3.4]{8.3.4}),
以及 $F_i$ 的各族 $X\mapsto\Cov(X)$ 在基变换下稳定这一事实.
证明第一个断言.
显然, 各 $\Cov(Y)$, $Y\in\underrightarrow{F}$,
中的族关于使函子 $\pi$ 连续的最粗拓扑 $T$ 覆盖 (III 1).
因此, 这些族生成的拓扑 $T'$ 比 $T$ 粗.
要证明它比\pageoriginale{} $T$ 细 (从而等于 $T$),
只须证明 $T'$ 的每个层 $M$ 都使 $M\circ\pi$ 为 $F$ 上的层;
或者等价地 (\hyperref[VI.rem8.2.4]{8.2.4}),
对每个 $i\in\ob I$, $M\circ\pi_i$ 为 $F_i$ 上的层.
由于函子 $\pi_i:F_i\to\underrightarrow{F}$ 与纤维积交换
(\hyperref[VI.lem8.3.4]{8.3.4}),
$M$ 为 $T'$ 的层, 当且仅当
(II \hyperref[II.cor2.3]{2.3} 及 I \hyperref[I.prop2.12]{2.12})
对每个 $i\in\ob I$、每个 $Y\in\ob F_i$ 及每个
$(Y_{\alpha}\to Y)_{\alpha\in A}\in\Cov(Y)$,
集合序列
$$
M(\pi_i(Y))\to\prod_{\alpha\in A}M(\pi_i(Y_{\alpha}))
\rightrightarrows\prod_{(\alpha,\beta)\in A\times A}
M\bigl(\pi_i(\fprod{Y_{\alpha}}{Y_{\beta}}{Y})\bigr)
$$
正合,
\ie (\loccit) 当且仅当对每个 $i\in\ob I$,
$M\circ\pi_i$ 为 $F_i$ 上的层.


% END INCLUDED SGA4_English_current_buildable_source_20260906/source/06/drafts/SGA4_VI_LINES_5857_6042_English.texfrag


% BEGIN INCLUDED SGA4_English_current_buildable_source_20260906/source/06/drafts/SGA4_VI_LINES_6043_6336_English.texfrag
\begin{enonce*}{命题 8.3.6}\etiquette[8.3.6]{VI.prop8.3.6}
沿用 \hyperref[VI.subsubsec8.3.1]{8.3.1}
及 \hyperref[VI.subsubsec8.3.2]{8.3.2} 的记号与假设.
若对每个 $i\in\ob I$,
$X\mapsto\Cov(X)$ 是 $F_i$ 上的预拓扑
\emph{(II \hyperref[II.def1.3]{1.3})},
且对每个 $X\in\ob F_i$, $\Cov(X)$ 中的覆盖族有限,
则对每个 $Y\in\underrightarrow{F}$,
$\Cov(Y)$ 中的覆盖族有限,
而 $Y\mapsto\Cov(Y)$ 为 $\underrightarrow{F}$ 上的预拓扑.
\end{enonce*}

\setcounter{subsubsection}{6}
\subsubsection{}\etiquette[8.3.7]{VI.subsubsec8.3.7}
只须证明各族 $Y\mapsto\Cov(Y)$
具有 (II \hyperref[II.def1.3]{1.3}) 的性质 (PT2)
(复合下稳定性).
设 $i\in\ob I$, $Y$ 是 $F_i$ 的对象,
$(Y_\alpha\to Y)_{\alpha\in A}\in\Cov(Y)$.
再设 $i_\alpha$, $\alpha\in A$, 是 $I$ 的一族对象;
对每个 $\alpha$, 设 $Z_\alpha$ 是 $F_i$ 的对象,
$(Z_{\alpha\beta}\xrightarrow{n_{\alpha\beta}}Z_\alpha)_{\alpha\beta\in
B_\alpha}$ 为 $\Cov(Z_\alpha)$ 中的族.
最后, 对每个 $\alpha\in A$, 给定一个同构
$\phi:\pi(Z_\alpha)\xrightarrow{\sim}\pi(Y_\alpha)$.
回到各 $\Cov(X)$, $X\in\ob\underrightarrow{F}$,
中的族的定义 (\hyperref[VI.subsubsec8.3.2]{8.3.2}),
可见须证明族
$$
(\pi (Z_{\alpha \beta}) \xrightarrow{\pi(n_{\alpha}) \circ \phi_{\alpha} \circ
  \pi(n_{\alpha \beta})} \pi(Y))_{\alpha \beta \in \mathop{\coprod}_A
  B_{\alpha}}
$$
属于 $\Cov(\pi(Y))$.
我们经过五次归结来证明这一点.

\setcounter{subsubsection}{7}
\subsubsection{}\etiquette[8.3.8]{VI.subsubsec8.3.8}
\emph{归结为对每个 $\alpha\in A$ 都有}
$\phi_\alpha=\pi(m_\alpha)$,
$m_\alpha:Z_\alpha\to Y_\alpha$ \emph{的情形}.
\emph{对每个 $\alpha\in A$, 有
$\phi_\alpha=\pi(m_\alpha)\pi(s_\alpha)^{-1}$
$(\hyperref[VI.prop6.5]{6.5})$,
其中 $s_\alpha$ 是 $f_\alpha:i'_\alpha\to i_\alpha$ 上的笛卡儿态射.
因此有交换图}
\[
\xymatrix@C=2.5cm@R=1.25cm{
Z_{\alpha\beta} \ar[d]_{n_{\alpha\beta}} & f^{\ast}(Z_{\alpha\beta})
\ar[l]_{s_{\alpha\beta}} \ar[d]^{f^{\ast}_{\alpha}(n_{\alpha\beta})}\\
Z_{\alpha} &   f^{\ast}_{\alpha}(Z_{\alpha}) \ar[l]_{s_{\alpha}},  &
**[l]\alpha\beta\in \displaystyle{\coprod_A} B_{\alpha},
}
\]
它们\pageoriginale{}给出交换图
\[
\xymatrix@C=2.5cm@R=1.25cm@H=.2cm{
\pi(Z_{\alpha\beta})  \ar[r]^{\sim} \ar[d]^{\pi(n_{\alpha\beta})}&
\pi(f^{\ast}_{\alpha}(Z_{\alpha\beta})) \ar[d]^{\pi(f^{\ast}_{\alpha}(m_{\alpha\beta}))}\\
\pi(Z_{\alpha})  \ar[r]^{\displaystyle\mathop{\sim}^{\pi(S_\alpha)^{-1}}} & \pi(f^{\ast}_{\alpha}(Z_{\alpha}))
\ar[r]^{\pi(m_{\alpha})} & \pi(Y_\alpha),\quad
\alpha\beta\in \displaystyle{\coprod_A} B_{\alpha}.}
\]
由于各族
$(f^{\ast}_{\alpha}(Z_{\alpha\beta})\xrightarrow{f^{\ast}(n_{\alpha})}
f^{\ast}(Z_{\alpha}))$
属于 $\Cov(f^{\ast}_{\alpha}(Z_{\alpha}))$,
可归结为对每个 $\alpha$ 都有 $\phi_\alpha=\pi(m_\alpha)$ 的情形.

\subsubsection{}\etiquette[8.3.9]{VI.subsubsec8.3.9}

\emph{归结为对每个 $\alpha\in A$, 有 $i_\alpha=j$ 且}
$p(m_\alpha)=k:j\to i$ \emph{的情形}.
置 $g_\alpha=p(m_\alpha)$.
由于 $I$ 余滤过且 $A$ 有限,
存在 $I$ 的一个对象 $j$ 及态射 $h_\alpha:j\to i_\alpha$,
使对每对 $(\alpha,\alpha')$,
有 $g_\alpha h_\alpha=g_{\alpha'}h_{\alpha'}=k$.
如 \hyperref[VI.subsubsec8.3.8]{8.3.8}
使用逆像函子 $h_\alpha^{\ast}$, 就归结为所述情形.

\subsubsection{}\etiquette[8.3.10]{VI.subsubsec8.3.10}

\emph{归结为 $i=j$ 且 $k:j\to i$ 为恒等的情形}.
对每个 $\alpha$, 有交换图
\[
\xymatrix@C=2.5cm@R=1.5cm{
Z_{\alpha}  \ar[r]^{m'_\alpha} &  k^{\ast}(Y_{\alpha}) \ar[d]_{k^{\ast}(n_{\alpha})}
\ar[r]^{t_{\alpha}} & Y_{\alpha} \ar[d]^{n_{\alpha}}\\
                    & k^{\ast}(Y)  \ar[r]^{t} & Y,
}
\]
其中态射 $t$ 及 $t_\alpha$, $\alpha\in A$, 为笛卡儿的,
$m'_\alpha$ 位于 $j$ 的恒等之上,
且 $t_\alpha\circ m'_\alpha=m_\alpha$.
将族 $(Y_\alpha\to Y)\alpha\in A$
换为\pageoriginale{}族
$(k^{\ast}(Y_\alpha)\to k^{\ast}(Y))_{\alpha\in A}\in
\Cov(k^{\ast}(Y_\alpha))$,
并将态射 $m_\alpha$ 换为态射 $m'_\alpha$,
就归结为所述情形.

\begin{enonce*}{引理 8.3.11}\etiquette[8.3.11]{VI.lem8.3.11}
设 $j$ 是 $I$ 的对象,
$m:X'\to X$ 是 $F_j$ 中使 $\pi(m)$ 为同构的态射.
存在态射 $\ell:j'\to j$,
使 $\ell^{\ast}(m)$ 为同构.
\end{enonce*}

证明存在态射 $f:i\to j$
及态射 $q:f^{\ast}(X)\to f^{\ast}(X')$,
使 $f^{\ast}(m)q=\id_{f^{\ast}(X)}$,
且 $\pi(q)$ 为同构.
事实上, 存在同构 $n:\pi(X)\to\pi(X')$,
使 $\pi(m)\circ n=\id_{\pi(X)}$.
因此存在 (\hyperref[VI.prop6.5]{6.5}) 两个态射
$X\xleftarrow{s_1}Y_1\xrightarrow{q_1}X'$,
其中 $s_1$ 为笛卡儿的,
使 $n=\pi(q_1)\pi(s_1)^{-1}$.
因此 $\pi(m)\circ\pi(q_1)=\pi(s_1)$.
从 $Y_1$ 到 $X$ 的态射 $mq_1$ 及 $s_1$
在 $\underrightarrow{F}$ 中有相同的像.
因此 (\cf 归纳极限中态射的描述
(\hyperref[VI.prop6.5]{6.5})),
存在笛卡儿态射 $s_2:Y_2\to Y_1$,
使 $mq_1s_2=s_1s_2$.
置 $q_1s_2=q_2$,
并以 $s_3$ 表示笛卡儿态射 $s_1s_2$.
于是 $mq_2=s_3$.
态射 $q_2:Y_2\to X'$ 本质唯一地分解为
$Y_2\xrightarrow{q}Y_3\xrightarrow{s_4}X'$,
其中 $s_4$ 为笛卡儿的,
$q$ 位于 $p(Y_2)$ 的恒等之上.
因此有交换图
\[\vcenter{
\xymatrix@C=2.25cm@R=1.5cm{
Y_3 \ar[d]_{s_4} & Y_2 \ar[l]_{q}\ar[d]^{s_3}\\
X'  \ar[r]^{m} & X.          .
}}\leqno{(8.3.11.1)}
\]\etiquette[8.3.11.1]{VI.eq8.3.11.1}
置 $p(s_4)=p(s_3)=f:i\to j$.
有 $f^{\ast}(X')\simeq Y_3$, $f^{\ast}(X')\simeq Y_2$,
且 $f^{\ast}(m):Y_2\to Y_3$ 使下图交换:
\[
\xymatrix@C=2.25cm@R=1.5cm{
Y_3  \ar[r]^{f^{\ast}(m)} \ar[d]_{s_4} & Y_2 \ar[d]^{s_3}\\
X'  \ar[r]^{m} & X.
}
\]
于是\pageoriginale{} $s_3=s_3(f^{\ast}(m)\circ q)$;
由于 $s_3$ 为笛卡儿的,
且 $f^{\ast}(m)\circ q$ 位于 $p(Y_2)$ 的恒等之上,
有 $f^{\ast}(m)\circ q=\id_{Y_2}$.
此外, 图 \eqhyperref[VI.eq8.3.11.1]{8.3.11.1} 的交换性表明
$\pi(q)$ 为同构.

再将前述结果应用于态射
$q:f^{\ast}(X)\to f^{\ast}(X')$:
存在态射 $f':j'\to i$
及态射 $q':f'^{\ast}(X')\to f'^{\ast}(X)$,
使 $f'^{\ast}(q)q'=\id_{f'^{\ast}(f'^{\ast}(X'))}$.
但置 $\ell^{\ast}=ff'$, 则有同构
$f'^{\ast}f^{\ast}\simeq\ell^{\ast}$.
因此有态射 $f'^{\ast}(q):\ell^{\ast}(X)\to\ell^{\ast}(X')$
及同构 $f'^{\ast}f^{\ast}\simeq\ell^{\ast}$.
因此有态射 $f'^{\ast}(q):\ell^{\ast}(X)\to\ell^{\ast}(X')$
及态射 $q':\ell^{\ast}(X')\to\ell^{\ast}(X)$,
使 $\ell^{\ast}(m){f'}^{\ast}(q)=\id_{\ell^{\ast}(X)}$
且 ${f'}^{\ast}(q)q'=\id_{\ell^{\ast}(X')}$.
故 $\ell^{\ast}(m)$ 为同构.

\setcounter{subsubsection}{11}
\subsubsection{}\etiquette[8.3.12]{VI.subsubsec8.3.12}
\emph{证明的结束}.
由于 $I$ 余滤过且 $A$ 有限,
由 \hyperref[VI.lem8.3.11]{8.3.11} 可知,
存在一个将所有态射 $m_\alpha$, $\alpha\in A$,
映为同构的逆像函子 $\ell^{\ast}$.
于是归结为各 $m_\alpha$ 都是同构的情形.
但此时覆盖族
$(Z_{\alpha\beta}\to Z)_{\alpha\beta\in B_\alpha}$
由覆盖族
$(Y_{\alpha\beta}\to Y_\alpha)_{\alpha\beta\in B_\alpha}\in
\Cov(Y_\alpha)$
经基变换 $m_\alpha:Z_\alpha\to Y_\alpha$ 得到.
因此归结为 $Z_\alpha=Y_\alpha$
且 $\Phi_\alpha=\id_{\pi(Y_\alpha)}$ 的情形.
但此时族
$(Y_{\alpha\beta}\xrightarrow{n_\alpha\circ n_{\alpha\beta}}Y)_{
\alpha\beta\in\mathop{\coprod}_A B_\alpha}$
属于 $\Cov(Y)$ (性质 (PT2)).
故其在 $\pi$ 下的像属于 $\Cov(\pi(Y))$
(\hyperref[VI.subsubsec8.3.2]{8.3.2}).

\begin{enonce*}{定理 8.3.13}\etiquette[8.3.13]{VI.thm8.3.13}
设 $p:F\to I$ 是一个在本质小余滤过范畴上纤维化的 $𝒰$-拓扑斯.
若对 $I$ 的每个对象 $i$,
纤维化拓扑斯 $F_i$ 凝聚
$(\hyperref[VI.def2.3]{2.3})$,
且对每个态射 $f:i\to j$,
拓扑斯态射 $f\bigdot=(f^{\ast},f_{\ast}):F_i\to F_j$
凝聚 (\hyperref[VI.def3.1]{3.1}),
则拓扑斯 $\limtop{I}F_i$ 凝聚,
且对 $I$ 的每个对象 $i$,
典范拓扑斯态射
$$
\mu_i : \limtop{I} F_i \to F_i
$$
凝聚.

此外, 若 $F_{\coh}$ 表示 $F$ 的满子范畴,
由在其纤维中凝聚的 $F$ 的对象
(\hyperref[VI.def1.13]{1.13}) 定义,
则范畴 $F_{\coh}$ 在 $I$ 上纤维化,
而 $\underrightarrow{F}_{\coh}$
典范等价于 $\limtop{I}F_i$ 的凝聚对象范畴.
\end{enonce*}

对每个\pageoriginale{} $i\in\ob I$,
以 $(F_i)_{\coh}$ 表示 $F_i$ 的凝聚对象
(\hyperref[VI.def1.13]{1.13}) 构成的 $𝒰$-景
(配备诱导拓扑).
由于 $F_i$ 凝聚, $F_i$ 为 $(F_i)_{\coh}$ 上的层拓扑斯.
对每个态射 $f:i\to j$,
态射 $f\bigdot=(f^{\ast},f_{\ast})$ 凝聚,
故 $f^{\ast}(F_j)_{\coh}\subset(F_i)_{\coh}$.
因此存在纤维化 $𝒰$-景 $p_{\coh}:F_{\coh}\to I$,
其纤维化景为各景 $(F_i)_{\coh}$,
其相伴纤维化拓扑斯等价于 $p:F\to I$.
故 (\hyperref[VI.thm8.2.3]{8.2.3}),
\begin{equation*}
\limtop{I, f.} F_i \approxeq
(\underrightarrow{F}_{\coh})^{\sim}\tag{8.3.13.1}
\end{equation*}\etiquette[8.3.13.1]{VI.eq8.3.13.1}
由于各拓扑斯 $F_i$ 凝聚,
$(F_i)_{\coh}$ 中有限积与纤维积可表示
(\hyperref[VI.prop2.2]{2.2}).
此外, 函子 $f^{\ast}:(F_i)_{\coh}\to(F_j)_{\coh}$ 左正合.
最后, 由于 $(F_i)_{\coh}$ 的对象拟紧,
有限覆盖族形成一个预拓扑.
于是由 (\hyperref[VI.lem8.3.4]{8.3.4}) 可知,
$\varprojlim_{I^{\circ},F}(F_i)_{\coh}$ 中有限积与纤维积可表示
(引理 \hyperref[VI.lem8.3.4]{8.3.4} 的补充:
证明典范函子 $(F_i)_{\coh}\to\underrightarrow{F}_{\coh}$
将终对象映为终对象);
由 (\hyperref[VI.prop8.3.6]{8.3.6}) 可知,
$\underrightarrow{F}_{\coh}$ 的对象拟紧.
故 (\hyperref[VI.subsubsec2.4.5]{2.4.5}),
$\limtop{I}(F_i)$ 凝聚.
最后, 典范态射 $\mu_i:\limtop{I}(F_i)\to F_i$
由景态射 (\hyperref[VI.thm8.2.3]{8.2.3})
$$
(F_i)_{\coh} \to \underrightarrow{F}_{\coh}
$$
通过过渡到层拓扑斯得到,
故 (\hyperref[VI.cor3.3]{3.3}) 各态射 $\mu_i$ 凝聚.

证明最后一个断言.
有 \eqhyperref[VI.eq8.3.13.1]{8.3.13.1}
$$
\limtop{ℐ} (F_i) \simeq
(\underrightarrow{F}_{\coh})^{\sim}{\quoi}.
$$
设 $X$ 是 $(F_{\coh})^{\sim}$ 的凝聚对象.
由于 $X$ 拟紧, 存在有限覆盖族 $(Y_\alpha\to X)$,
其中 $Y_\alpha\in\ob\underrightarrow{F}_{\coh}$
(\hyperref[VI.def1.1]{1.1}).
取指标 $i\in\ob I$ 充分小,
可假定各 $Y_\alpha$ 是一个族 $Y'_\alpha$
在 $\mu_i:F_{i\coh}\to\underrightarrow{F}_{\coh}$ 下的像;
于是取各 $Y'_\alpha$ 的直和
(\hyperref[VI.cor1.15]{1.15}),
可见存在指标 $i\in\ob I$、对象 $Y'\in F_{i\coh}$
及满射态射 $\mu_i(Y')\to X$.
等价关系 $R=\fprod{\mu_i(Y')}{\mu_i(Y')}{X}$
于是为凝聚对象\pageoriginale{},
因为 $(\underrightarrow{F}_{\coh})^{\sim}$ 是凝聚拓扑斯
(\hyperref[VI.prop2.2]{2.2}).
先证明 $R$ 是 $\underrightarrow{F}_{\coh}$ 的对象.
由定义, $R$ 是 $\mu_i(Y'\times Y')$ 的子对象,
而且由于 $R$ 凝聚,
在改变指标 $i$ 后可找到 $F_{i\coh}$ 的一个箭头
$m:Z\to(Y'\times Y')$,
使 $\text{Im}(\mu_i(m))=R$.
由于 $F_i$ 凝聚, $\text{Im}(m)$ 在 $F_i$ 中凝聚
(\hyperref[VI.cor1.17.1]{1.17.1}),
故属于 $F_{i\coh}$.
因此存在指标 $i$ 及 $F_{i\coh}$ 的一个图
$\def\labelstyle{\scriptstyle} \xymatrix@C=1cm{
R \ar@<3pt>[r]^{p_1}\ar@<-3pt>[r]_{p_2} & Y'
}
$
使
$\def\labelstyle{\scriptstyle}
\xymatrix@C=1cm{
  \mu_i(R)\ar@<3pt>[r]^{\mu_i(p_1)} \ar@<-3pt>[r]_{\mu_i(p_2)}& \mu_i(Y')
}
$
为等价关系, 且 $X=\mu_i(Y')/\mu_i(R)$.
置 $X_i=\coker(p_1,p_2)$
及 $R'=\fprod{Y'}{Y'}{X'_i}$.
$F_i$ 中存在典范态射 $u:R\to R'$.
此外, $\mu_i^{\ast}(u)$ 为同构.
因此, 改变指标 $i$ 后可假定 $R=R'$,
\ie $R$ 为等价关系.
但此时 $X_i$ 属于 $F_{i\coh}$
(\hyperref[VI.cor1.17.1]{1.17.1}),
故 $X=\mu_i(X_i)$ 属于 $\underrightarrow{F}_{\coh}$.

\begin{enonce*}{推论 8.3.14}\etiquette[8.3.14]{VI.cor8.3.14}
设 $I$ 是一个本质小余滤过范畴,
$p:F\to I$ 与 $q:G\to I$ 为 $I$ 上的两个纤维化 $𝒰$-拓扑斯,
使对 $I$ 的每个对象 $i$,
拓扑斯 $F_i$ 与 $G_i$ 凝聚,
且对每个态射 $f:i\to j$,
态射 $f.:F_i\to F_j$ 及 $G_i\to G_j$ 凝聚.
再设 $m:F\to G$ 为纤维化 $𝒰$-拓扑斯的笛卡儿态射,
使对每个 $i\in\ob(I)$,
$m_i:F_i\to G_i$ 凝聚.
则由 $m$ 过渡到投射极限所得的态射 $\underleftarrow{m}$ 凝聚.
\end{enonce*}

对 $I$ 的每个对象 $i$,
态射 $m_i$ 诱导函子
$m^{\ast}_{i,\coh}:G_{i,\coh}\to F_{i,\coh}$,
从而有笛卡儿函子
$m^{\ast}_{\coh}=G_{\coh}\to F_{\coh}$,
它为从纤维化景 $F_{\coh}$ 到纤维化景 $G_{\coh}$ 的态射
\hyperref[VI.thm7.4.13.1]{(7.4.13)}.
函子 $\varinjlim_{I^{\circ}}m^{\ast}_{\coh}:
\varinjlim_{I^{\circ}}G_{\coh}\to\varinjlim_{I^{\circ}}F_{\coh}$
为景态射, 过渡到相应拓扑斯后给出态射 $\underleftarrow{m}$.
断言因而来自 \hyperref[VI.thm8.3.13]{8.3.13}
及 \hyperref[VI.cor3.3]{3.3}.

\begin{enonce*}[remark]{习题 8.3.15}\etiquette[8.3.15]{VI.exe8.3.15}
沿用 \hyperref[VI.thm8.3.13]{8.3.13} 的记号,
证明若各 $F_i$, $i\in I$, 为代数的
(\hyperref[VI.def2.3]{2.3}),
且各态射 $f.:F_i\to F_j$, $f\in\text{\rm Fl}(I)$,
凝聚,\pageoriginale{}
则 $\limtop{I}F_i$ 为代数的,
且 $(\limtop{I}F_i)_{\coh}\xrightarrow{\sim}\underrightarrow{F}_{\coh}$.
\end{enonce*}


% END INCLUDED SGA4_English_current_buildable_source_20260906/source/06/drafts/SGA4_VI_LINES_6043_6336_English.texfrag


% BEGIN INCLUDED SGA4_English_current_buildable_source_20260906/source/06/drafts/SGA4_VI_LINES_6337_6444_English.texfrag
\subsection{例: 局部拓扑斯}\etiquette[8.4]{VI.subsec8.4}

\subsubsection{}\etiquette[8.4.1]{VI.subsubsec8.4.1}

设 $X=(X_i)_{i\in I}$ 是范畴 $I$ 上的纤维化拓扑斯.
对 $I$ 的每个 pro-对象 $\sheaf{\alpha}=(i_{\alpha})$,
用公式
\begin{equation*}
X_{\sheaf{\alpha}}=\limtop{I_{\alpha}}X_{i\alpha}\tag{8.4.1.1}
\end{equation*}\etiquette[8.4.1.1]{VI.eq8.4.1.1}
由此得到范畴 $\Pro(I)$ (I \hyperref[I.sec8.10]{8.10}) 上的纤维化拓扑斯
(为此, 可将习题 I \hyperref[VI.subsubsec8.2.8]{8.2.8} 推广到伪函子).

\subsubsection{}\etiquette[8.4.2]{VI.subsubsec8.4.2}

设 $E$ 是一个拓扑斯,
$p:\text{Fl}(E)\to E$ 是 \hyperref[VI.exem7.3.1]{7.3.1}
中考虑的纤维化拓扑斯.
如 \hyperref[VI.subsubsec8.4.1]{8.4.1} 将其扩张,
得到纤维化拓扑斯 $\bar{\text{Fl}(E)}\to\Pro(E)$.
范畴 $\Point(E)$ 由一个全忠实函子映入 $\Pro(E)$
(IV \hyperref[IV.subsubsec6.8.5]{6.8.5}),
从而经基变换 (\hyperref[VI.subsubsec7.1.9]{7.1.9})
得到纤维化拓扑斯 $\Loc(E)\to\Point(E)$.
在 $E$ 的点 $p$ 处的纤维 $\Loc_p(E)$,
称为\emph{$E$ 在点 $p$ 处的局部化}.
由前面所述, 此拓扑斯协变地依赖于点 $p$.
由定义,
\begin{equation*}
\Loc_p(E)=\Limtop_{X\in\Vois(p)}E_{/X}{\quoi},\tag{8.4.2.1}
\end{equation*}\etiquette[8.4.2.1]{VI.eq8.4.2.1}
其中 $\Vois(p)$ 是 $p$ 的邻域范畴
(IV \hyperref[IV.subsec6.8]{6.8}).

\subsubsection{}\etiquette[8.4.3]{VI.subsubsec8.4.3}

设 $m:p'\to p$ 是 $\Point(E)$ 的态射 (IV 6).
对每个 $X\in\Vois(p)$,
由此得到 $E_{/X}$ 的一个点 $p'_X$,
从而由 $\limtop{}$ 的泛性质 \hyperref[VI.subsec8.1]{8.1},
得到 $\Loc_p(E)$ 的一个点 $\theta_p(m)$.
这定义函子
\begin{equation*}
\theta_p:\Point(E)_{/p}\to\Point(\Loc_p(E)){\quoi},\tag{8.4.3.1}
\end{equation*}\etiquette[8.4.3.1]{VI.eq8.4.3.1}
立即可见它为范畴等价.
因此, $\Point(\Loc_p(E))$
典范等价于 $p$ 的一般化范畴
(IV \hyperref[IV.subsubsec7.1.8]{7.1.8}).
此外, 有拓扑斯的典范等价
$$
\can\bigdot:\Loc_{\theta_p(m)}(\Loc_p(E))\to\Loc_{p'}(E){\quoi},
$$
它置于在同构意义下交换的图中
\[
\xymatrix@R=1.5cm{
\Loc_{\theta_p(m)}(\Loc_p(E))\ar[rr]^{\can\bigdot}\ar[rd]_{\phi}
&&\Loc_{p'}(E)\ar[ld]^{m}\\
&\Loc_p(E)&
}
\]
其中 $\phi$\pageoriginale{}
为 \eqhyperref[VI.eq8.4.2.1]{8.4.2.1} 中的典范态射.

\subsubsection{}\etiquette[8.4.4]{VI.subsubsec8.4.4}

局部化拓扑斯似乎只对来自代数几何的拓扑斯,
或至少具有适当有限性性质的拓扑斯才有意义.
例如, 当 $E$ 是分离拓扑空间上的层拓扑斯时,
立即可见 (例如利用 \hyperref[VI.thm8.2.9]{8.2.9}),
局部化拓扑斯均为点拓扑斯.
此外, 在此情形范畴 $\Point(E)$ 离散,
\hyperref[VI.subsubsec8.4.2]{8.4.2} 描述的情形是平凡的.
另一方面, 设 $E=\Top(X)$ 是概形 $X$ 的 Zariski 拓扑的层拓扑斯,
$x$ 是 $X$ 的点,
$O_x$ 是 $X$ 在 $x$ 处的局部环,
且 $Y=\spec(O_x)$.
得到
$$
\Loc_x(E)=\Top(Y){\quoi}.
$$
来自平展拓扑的其他例子, 见本讨论班第 VIII 篇.

\subsubsection{}\etiquette[8.4.5]{VI.subsubsec8.4.5}

当 $E$ 局部凝聚 (\hyperref[VI.def2.3]{2.3}) 时,
局部化拓扑斯 $\Loc_p(E)$ 凝聚
(在 \eqhyperref[VI.eq8.4.2.1]{8.4.2.1} 中,
可以只取代数且凝聚的 $X\in\Vois(p)$,
然后应用 \hyperref[VI.thm8.3.13]{8.3.13}).
对每个点的态射 $m:p'\to p$,
拓扑斯态射 $m\bigdot:\Loc_{p'}(E)\to\Loc_p(E)$ 凝聚
(\hyperref[VI.cor8.3.14]{8.3.14}).

\subsubsection{}\etiquette[8.4.6]{VI.subsubsec8.4.6}

\emph{局部拓扑斯（local topos）}是使函子
$\Gamma(X,-)=$ “$X$ 上的截面”
(IV \hyperref[IV.subsec4.3]{4.3})
为纤维函子 (IV 6) 的拓扑斯 $X$.
对应于此纤维函子的点称为局部拓扑斯的\emph{中心}.
当 $E$ 局部凝聚时,
其局部化拓扑斯为局部拓扑斯
(\hyperref[VI.thm8.2.3]{1.2.3}、
\hyperref[VI.prop8.5.2]{8.5.2}
及 \hyperref[VI.cor8.5.7]{8.5.7}).
$\Loc_p(E)$ 的中心典范同构于 $\theta_p(\id_p)$
(\eqhyperref[VI.eq8.4.3.1]{8.4.3.1}).
$\Loc_p(E)$ 的中心在 $E$ 中的像同构于 $p$.
局部化拓扑斯 $\Loc_p(E)$,
配备典范态射\pageoriginale{} $\Loc_p(E)\to E$ 后,
是将局部拓扑斯映入 $E$,
使中心映“到” $p$ 这一泛 (2-泛!) 问题的解.

\subsubsection{}\etiquette[8.4.7]{VI.subsubsec8.4.7}

编者尚未处理关于局部拓扑斯的若干问题.
例如, 若 $X$ 是局部拓扑斯,
$X$ 的终对象 $e$ 有一个极大的开子拓扑斯 $U\ne e$.
$U$ 的补 $Y$ 是没有非平凡开子拓扑斯的局部拓扑斯.
这样的拓扑斯是点拓扑斯吗?
设 $U$ 为拓扑斯,
$X=(U,\Phi:U\to(\Ens))$ 是经粘合得到的拓扑斯.
粘合函子 $\Phi$ 应满足什么条件, 才能使 $X$ 为局部拓扑斯?


% END INCLUDED SGA4_English_current_buildable_source_20260906/source/06/drafts/SGA4_VI_LINES_6337_6444_English.texfrag


% BEGIN INCLUDED SGA4_English_current_buildable_source_20260906/source/06/drafts/SGA4_VI_LINES_6445_6610_English.texfrag
\subsection{拓扑斯的滤过投射极限的层的结构}\etiquette[8.5]{VI.subsec8.5}

\subsubsection{}\etiquette[8.5.1]{VI.subsubsec8.5.1}

本节中, $F=(F_i)_{i\in I}$ 是小余滤过范畴 $I$ 上的纤维化 $\Uu$-拓扑斯.
选定 $F$ 的一个双劈裂（bicleavage）
(\hyperref[VI.subsubsec7.1.4]{7.1.4}),
\ie 对每个 $f\in\text{Fl}(I):i\to j$,
选定拓扑斯态射 $f\bigdot:F_i\to F_j$,
使 $f^{\ast}:F_j\to F_i$ 为纤维化结构的逆像函子.
于是有满足余循环条件的典范同构 $c_{f,g}$
(\hyperref[VI.subsubsec7.1.3]{7.1.3}).
以 $\underleftarrow{F}$ 表示纤维化拓扑斯 $F$ 的投射极限
(\hyperref[VI.thm8.2.3]{8.2.3});
并对每个 $i\in\ob I$, 以
\begin{equation*}
\mu_i:\underleftarrow{F}\to F_i{\quoi},\tag{8.5.1.1}
\end{equation*}\etiquette[8.5.1.1]{VI.eq8.5.1.1}
表示典范态射 (\hyperref[VI.subsubsec8.1.3]{8.1.3}).
以 $\Top(F)$ 表示 $F$ 的总拓扑斯
\hyperlink{sga4.exp.VI.remark.7.4.3.item.3}{(7.4.3.3)}.
由 \hyperref[VI.thm8.2.9]{8.2.9},
典范函子 $F\to\underrightarrow{F}$ 定义拓扑斯态射
\begin{equation*}
Q:\underleftarrow{F}\to\Top(F){\quoi}.\tag{8.5.1.2}
\end{equation*}\etiquette[8.5.1.2]{VI.eq8.5.1.2}
将拓扑斯 $\underleftarrow{F}$ 等同于
$\Hhom_{\Cart/I^{\circ}}(I^{\circ},F')$
(\hyperref[VI.thm8.2.9]{8.2.9}),
并将拓扑斯 $\Top(F)$ 等同于
$\Hhom_{I^{\circ}}(I^{\circ},F')$
(\hyperref[VI.prop7.4.7]{7.4.7}).
在这些等同下, $Q$ 为
$\Hhom_{\cart/I^{\circ}}(I^{\circ},F')$ 到
$\Hhom_{I^{\circ}}(I^{\circ},F')$ 的嵌入态射
(\hyperref[VI.thm8.2.9]{8.2.9}),
并且对 $\underleftarrow{F}$ 的每个对象 $M$, 有
\begin{equation*}
Q_{\ast}(M)=(i\longmapsto\mu_{i\ast}(M)){\quoi}.\tag{8.5.1.3}
\end{equation*}\etiquette[8.5.1.3]{VI.eq8.5.1.3}
由 \hyperlink{sga4.exp.VI.remark.7.4.3.item.4}{7.4.3.4},
对每个 $i\in\ob I$, 存在拓扑斯态射
\begin{equation*}
\alpha_i:F_i\to\Top(F){\quoi}.\tag{8.5.1.4}
\end{equation*}\etiquette[8.5.1.4]{VI.eq8.5.1.4}
图\pageoriginale{}
\[\vcenter{
\xymatrix@R=1.25cm{
{\underleftarrow{F}}\ar[rr]^{\mu_i}\ar[rd]_{Q}&&F_i\ar[ld]^{\alpha_i}\\
&\Top(F)&
}}\leqno{(8.5.1.5)}
\]\etiquette[8.5.1.5]{VI.eq8.5.1.5}
一般不交换 (即使在同构意义下也不交换).
但所涉及态射的定义给出典范同构
\begin{equation*}
\alpha_i^{\ast}Q_{\ast}\simeq\mu_{i\ast}.\tag{8.5.1.6}{\quoi}
\end{equation*}\etiquette[8.5.1.6]{VI.eq8.5.1.6}
当 $i$ 是 $I$ 的\emph{终对象}时,
$\alpha_i$ 是具有缩回 $\beta_i:\Top(F)\to F_i$ 的嵌入
(\hyperref[VI.lem7.4.12]{7.4.12}),
而图
\[\vcenter{
\xymatrix@R=1.25cm{
\underleftarrow{F}\ar[rr]^{\mu_i}\ar[rd]_{Q}&&F_i\ar[ld]^{\beta_i}\\
&\Top(F)&,
}}\leqno{(8.5.1.7)}
\]\etiquette[8.5.1.7]{VI.eq8.5.1.7}
在典范同构意义下交换.

\begin{enonce*}{命题 8.5.2}\etiquette[8.5.2]{VI.prop8.5.2}
设 $j\mapsto M_j$ 为 $\Top(F)$ 的一个对象.
有函子性同构
\begin{equation*}
Q^{\ast}(j\mapsto M_j)\simeq\mathop{\varinjlim}_{I^{\circ}}
\mu_j^{\ast}M_j{\quoi}.\tag{8.5.2.1}
\end{equation*}\etiquette[8.5.2.1]{VI.eq8.5.2.1}
\end{enonce*}
\noindent
$\Top(F)$ 的一个对象
(\hyperref[VI.subsubsec7.1.3]{7.1.3})
由如下数据构成:
一个对应 $j\mapsto M_j$, $M_j\in\ob F_j$,
以及对每个态射 $f:i\to j$, 一个态射
\begin{equation*}
\beta_f:M_j\to f_{\ast}M_i\tag{8.5.2.2}
\end{equation*}\etiquette[8.5.2.2]{VI.eq8.5.2.2}
或者由伴随性等价地, 一个态射
\begin{equation*}
\beta'_f:f^{\ast}M_j\to M_i{\quoi},\tag{8.5.2.3}
\end{equation*}\etiquette[8.5.2.3]{VI.eq8.5.2.3}
其中各态射 $\beta'_f$ 须满足条件:
对每对可复合态射 $i\xrightarrow{f}j\xrightarrow{g}k$,
下图交换:
\[\vcenter{
\xymatrix@C=2.25cm@R=1.5cm{
f^{\ast}g^{\ast}M_k\ar[r]^{f^{\ast}\beta'_g}\ar@{-}[d]_{c^{\ast}_{g,f}}&
f^{\ast}M_j\ar@{->}[d]^{\beta'_f}\\
(gf)^{\ast}M_k\ar[r]^{\beta'_{gf}}&M_i
}}\leqno{(8.5.2.4)}
\]\etiquette[8.5.2.4]{VI.eq8.5.2.4}
再回忆\pageoriginale{} \eqhyperref[VI.eq8.1.3.1]{8.1.3.1},
对每个态射 $f:i\to j$,
存在同构 $b_f:f\bigdot\mu_i\xrightarrow{\sim}\mu_j$,
并且对每对可复合态射 $i\xrightarrow{f}j\xrightarrow{g}k$,
有交换图
\[
\xymatrix@C=2cm@R=1.25cm{
g\bigdot f\bigdot\mu_i\ar[r]^{c_{g,f}}\ar[d]^{g\bigdot(b_f)}&
(gf)\bigdot\mu_i\ar[d]^{b_{gf}}\\
g\bigdot\mu_j\ar[r]^{b_g}&\mu_k
}\leqno{(8.5.2.5)}
\]\etiquette[8.5.2.5]{VI.eq8.5.2.5}
设 $f:i\to j$ 为 $I$ 的一个态射.
以
\begin{equation*}
t_f:\mu_j^{\ast}(M_j)\to\mu_i^{\ast}(M_i)\tag{8.5.2.6}
\end{equation*}\etiquette[8.5.2.6]{VI.eq8.5.2.6}
表示复合
$\mu_j^{\ast}(M_j)\xrightarrow{b_f^{\ast}}
\mu_i^{\ast}(f^{\ast}(M_j))
\xrightarrow{\mu_i^{\ast}(\beta'_f)}\mu_i^{\ast}(M_i)$.
请读者验证, 图 \eqhyperref[VI.eq8.5.2.4]{8.5.2.4}
及 \eqhyperref[VI.eq8.5.2.5]{8.5.2.5} 的交换性蕴含,
对每对可复合态射 $i\to j\to k$, 有
\begin{equation*}
t_ft_g=t_{gf}{\quoi}.\tag{8.5.2.7}
\end{equation*}\etiquette[8.5.2.7]{VI.eq8.5.2.7}
由此定义了函子 $I^{\circ}\to\underleftarrow{F}$,
可以考察其归纳极限 $\varinjlim_{I^{\circ}}\mu_i^{\ast}(M_i)$.
设 $N$ 为 $\underleftarrow{F}$ 的一个对象.
有
\begin{equation*}
\Hom(\varinjlim_{I^{\circ}}\mu_i^{\ast}(M_i),N)\simeq
\varprojlim_{I^{\circ}}\Hom(\mu_i^{\ast}(M_i),N){\quoi},\tag{8.5.2.8}
\end{equation*}\etiquette[8.5.2.8]{VI.eq8.5.2.8}
从而由伴随性有同构
\begin{equation*}
\Hom(\varinjlim_{I^{\circ}}\mu_i^{\ast}(M_i),N)\simeq
\varprojlim_{I^{\circ}}\Hom_{F_i}(M_i,\mu_{i\ast}(N)){\quoi}.
\tag{8.5.2.9}
\end{equation*}\etiquette[8.5.2.9]{VI.eq8.5.2.9}
现在\pageoriginale{}解释 \eqhyperref[VI.eq8.5.2.9]{8.5.2.9} 的右端.
对每个态射 $f:i\to j$, 以
$$
d_f:\Hom_{F_i}(M_i,\mu_{i\ast}(N))\to
\Hom_{F_j}(M_j,\mu_{j\ast}(N)){\quoi}
$$
表示 \eqhyperref[VI.eq8.5.2.9]{8.5.2.9}
中出现的投射系的转移映射.
由态射 $t_f$ 的定义
(\eqhyperref[VI.eq8.5.2.6]{8.5.2.6}) 可知,
$d_f$ 给态射
$$
u_i:M_i\to\mu_{i\ast}(N)
$$
对应由态射
$$
M_j\xrightarrow{\beta_f}f_{\ast}(M_i)
\xrightarrow{f_{\ast}(u_i)}f_{\ast}\mu_{i\ast}(N)
\xrightarrow{b_f}\mu_{j\ast}(N){\quoi}.
$$
复合得到的态射
$$
d_f(u_i):M_j\to\mu_{j\ast}(N){\quoi},
$$
因此, \eqhyperref[VI.eq8.5.2.9]{8.5.2.9} 右端的元素,
是一个态射族 $u_i:M_i\to\mu_{i\ast}(N)$, $i\in\ob I$,
使对每个 $f:i\to j$, 图
\[
\xymatrix@C=2cm@R=1.25cm{
M_j\ar[r]\ar[d]_{u_j}&f_{\ast}(M_i)\ar[d]^{f_{\ast}(u_i)}\\
\mu_{j\ast}(N)\ar[r]^{\sim}&f_{\ast}\mu_{i\ast}(N)
}
\]
交换, 即从截面
$(i\mapsto M_i)\in\Hhom_{I^{\circ}}(I^{\circ},F')$
到截面 $Q_{\ast}(N)$ 的一个态射.
因此有函子性同构
$$
\Hom(\varinjlim_{I^{\circ}}\mu_i^{\ast}(M_i),N)
\simeq\Hom((i\mapsto M_i),Q_{\ast}(N)){\quoi};
$$
命题由伴随性得到.


% END INCLUDED SGA4_English_current_buildable_source_20260906/source/06/drafts/SGA4_VI_LINES_6445_6610_English.texfrag


% BEGIN INCLUDED SGA4_English_current_buildable_source_20260906/source/06/drafts/SGA4_VI_LINES_6611_6839_English.texfrag
\begin{enonce*}{命题 8.5.3}\etiquette[8.5.3]{VI.prop8.5.3}
若对每个 $f:i\to j\in\text{\rm Fl}(I)$，函子
$f_{\ast}:F_i\to F_j$ 都与小滤过归纳极限交换，则对每个截面
$(i\mapsto M_i)\in\Hom_I^{\circ}(I^{\circ},F')$，有
\begin{equation*}
Q_{\ast} Q^{\ast} (i \mapsto M_i) \simeq i \mapsto
\varinjlim_{f : j \to i} f_{\ast} M_j{\quoi}.\tag{8.5.3.1}
\end{equation*}\etiquette[8.5.3.1]{VI.eq8.5.3.1}
\end{enonce*}

\noindent 我们\pageoriginale{}将利用 $\beta_f$ \eqhyperref[VI.eq8.5.2.2]{8.5.2.2} 及
$c_{f,g}$（\hyperref[VI.subsubsec8.5.1]{8.5.1}）显式描述归纳系
$(f:j\to i)\mapsto f_{\ast}(M_j)$ 的过渡态射以及截面
$i\mapsto\varinjlim_{f:j\to i}f_{\ast}M_j$ 的过渡态射留给读者。由于函子
$f_{\ast}$ 与滤过归纳极限交换，截面
$i\mapsto\varinjlim_{f:j\to i}f_{\ast}M_j$ 是笛卡儿的。此外，从
$(i\mapsto M_i)$ 到截面
$(i\mapsto\varinjlim_{f:j\to i}f_{\ast}M_j)$ 有一个自然态射 $u$；显然，从
$(i\mapsto M_i)$ 到任意笛卡儿截面的每个态射都唯一地通过 $u$ 分解，由此得到同构
\eqhyperref[VI.eq8.5.3.1]{8.5.3.1}。

\begin{enonce*}{推论 8.5.4}\etiquette[8.5.4]{VI.cor8.5.4}
在 \hyperref[VI.prop8.5.3]{8.5.3} 的假设下，函子 $Q_{\ast}$ 与小滤过归纳极限交换。
\end{enonce*}

设 $\alpha\mapsto N^{\alpha}$ 是 $\underleftarrow{F}$ 中的一个小滤过归纳系。令
$M^{\alpha}=Q_{\ast}(N^{\alpha})$，于是 $M^{\alpha}$ 是截面
$i\mapsto M_i^{\alpha}$。有 $N^{\alpha}\simeq Q^{\ast}Q_{\ast}(N^{\alpha})$；由于
$Q^{\ast}$ 与归纳极限交换，有
$\varinjlim_{\alpha}N^{\alpha}\simeq Q^{\ast}(\varinjlim_{\alpha}M^{\alpha})$。
$\Top(F)=\Hom_{I^{\circ}}(I^{\circ},F')$ 中的归纳极限逐纤维计算。因此
$\varinjlim_{\alpha}M^{\alpha}\simeq (i\mapsto\varinjlim_{\alpha}M_i^{\alpha})$。
由 \eqhyperref[VI.eq8.5.3.1]{8.5.3.1}，有
$$
Q_{\ast} (\varinjlim_{\alpha} {N}^{\alpha}) \simeq Q_{\ast}
Q^{\ast} (\varinjlim_{\alpha} {M}^{\alpha}) \simeq i \mapsto
\varinjlim_{f : j \to i} f_{\ast}
(\varinjlim_{\alpha} {M}^{\alpha}_j){\quoi}.
$$
由于函子 $f_{\ast}$ 与滤过归纳极限交换，得到
$$
Q_{\ast} (\varinjlim_{\alpha} {N}^{\alpha}) \simeq i \mapsto
\varinjlim_{f: j\to i} \varinjlim_{\alpha} f_{\ast}
(M^{\alpha}_j) \simeq i \mapsto \varinjlim_{\alpha} \varinjlim_{f:
  j\to i} f_{\ast}
(M^{\alpha}_j) \simeq \varinjlim_{\alpha} Q_{\ast} (N^{\alpha}){\quoi}.
$$

\begin{enonce*}{推论 8.5.5}\etiquette[8.5.5]{VI.cor8.5.5}
若函子 $f_{\ast}:F_i\to F_j$ 与小滤过归纳极限交换，则对 $I$ 的每个对象 $i$ 及每个
$M_i\in\ob F_i$，有
\begin{equation*}
\mu_{i\ast} \mu_{i}^{\ast} (M_i) \simeq \varinjlim_{f : j \to
      i} f_{\ast} f^{\ast}(M_i){\quoi}.\tag{8.5.5.1}
\end{equation*}\etiquette[8.5.5.1]{VI.eq8.5.5.1}
\end{enonce*}
\noindent
作基变换 $I/i\to I$ 后，可以假设 $i$ 是 $I$ 的终对象（\hyperref[VI.lem8.2.1]{8.2.1}）。
此时态射 $\mu_i:\underleftarrow{F}\to F_i$ 是下列态射的复合\pageoriginale{}（\hyperref[VI.thm8.2.9]{8.2.9} 及
\hyperref[VI.lem7.4.12]{7.4.12}）：
\begin{align*}
Q &: \underleftarrow{F} \to F^{\sim}{\quoi},\\
\beta_i &: (\alpha_{i!}, \alpha^{\ast}_i) : F^{\sim} \to F_i{\quoi}.
\end{align*}
因此 $\mu_i^{\ast}(M_i)\simeq Q^{\ast}\beta_i^{\ast}(M_i)$，其中
$\beta_i^{\ast}(M_i)$ 是截面 $(f:i\to j)\mapsto f^{\ast}(M_i)$，其过渡态射由
$c_{f,g}$ 导出。所断言的结论于是由 \hyperref[VI.prop8.5.3]{8.5.3} 得到。

\begin{enonce*}{推论 8.5.6}\etiquette[8.5.6]{VI.cor8.5.6}
在 \hyperref[VI.cor8.5.5]{8.5.5} 的假设下，函子
$$
\mu_{i\ast} : \underleftarrow{F} \to F_i
$$
与滤过归纳极限交换。
\end{enonce*}

函子 $\mu_{i\ast}$ 是函子 $Q_{\ast}$ 与“限制到 $i$ 处纤维”这一函子的复合；前者与滤过归纳极限交换，后者与归纳极限交换。

\begin{enonce*}{推论 8.5.7}\etiquette[8.5.7]{VI.cor8.5.7}
在 \hyperref[VI.prop8.5.3]{8.5.3} 的假设下，设 $i$ 是 $I$ 的对象，$X$ 是 $F_i$ 的对象，且
$F_i$ 上的函子 $\Hom(X,-)$ 与滤过归纳极限交换。则函子
$\Hom(\mu_i^{\ast}(X),-)$ 与滤过归纳极限交换，并且对 $F_i$ 的每个对象 $M_i$，有
\begin{equation*}
\Hom(\mu^{\ast}_i(X), \mu^{\ast}_i(M_i)) \simeq \varinjlim_{f
: j \to i} \Hom(f^{\ast}(X), f^{\ast}(M_i))\tag{8.5.7.1}{\quoi}.
\end{equation*}\etiquette[8.5.7.1]{VI.eq8.5.7.1}
\end{enonce*}
\noindent
对 $\Top(F)$ 的每个对象 $(j\mapsto M_j)$，有
\begin{equation*}
\Hom(\mu^{\ast}_i(X), Q^{\ast}(j \mapsto M_j)) \simeq \varinjlim_{f
: j \to i} \Hom(f^{\ast}(X), M_{j}){\quoi}.\tag{8.5.7.2}
\end{equation*}\etiquette[8.5.7.2]{VI.eq8.5.7.2}

函子 $\Hom(\mu_i^{\ast}(X),-)$ 同构于函子 $\Hom(X,\mu_{i\ast}(-))$，后者与滤过归纳极限交换（\hyperref[VI.cor8.5.4]{8.5.4}）。
由 \eqhyperref[VI.eq8.5.5.1]{8.5.5.1}，有
$$
\Hom(\mu^{\ast}_i(X), \mu^{\ast}_i (M_i)) \simeq \Hom (X, \mu_{i\ast}\mu^{\ast}_i (M_i)) \simeq  \Hom(X, \varinjlim_{f : j \to
  i} f_{\ast} f^{\ast} (M_i)){\quoi},
$$
由此得到公式 \eqhyperref[VI.eq8.5.7.1]{8.5.7.1}。同样，由
\eqhyperref[VI.eq8.5.3.1]{8.5.3.1}，有
$$
\Hom(\mu^{\ast}_i (X), Q^{\ast}(j \mapsto M_j)) \simeq \Hom(X,
\mu_{i\ast} Q^{\ast} (j \mapsto M_j)) \simeq \Hom(X,
\varinjlim_{f : j \to i} f_{\ast} (M_j)){\quoi},
$$
由此\pageoriginale{}得到公式 \eqhyperref[VI.eq8.5.7.2]{8.5.7.2}。

\begin{enonce*}{命题 8.5.8}\etiquette[8.5.8]{VI.prop8.5.8}
设 $G\to I$ 是纤维化 $𝒰$-拓扑斯，$m:F\to G$ 是纤维化拓扑斯的笛卡儿态射
\hyperref[VI.def7.1.5]{(7.1.5)}。对 $G$ 使用 \hyperref[VI.def8.1.1]{8.1.1} 中引入的记号。
设 $\underleftarrow{m}:\underleftarrow{F}\to\underleftarrow{G}$ 是由 $m$ 取投射极限而得到的态射
（\hyperref[VI.subsubsec8.1.4]{8.1.4}）。下列拓扑斯及拓扑斯态射组成的图
\[\vcenter{
\xymatrix@C=2.25cm@R=1.25cm{
\underset{\leftarrow}{F}  \ar[r]^{Q} \ar[d]^{\underset{\leftarrow}{m}}& \Top(F)\ar[d]^{m^{\sim}}\\
\underset{\leftarrow}{G}  \ar[r]^{Q} & \Top(G)
}}\leqno{(8.5.8.1)}
\]\etiquette[8.5.8.1]{VI.eq8.5.8.1}
在同构意义下交换。对 $\underleftarrow{F}$ 的每个对象 $N$，有
\begin{equation*}
\underleftarrow{m}_{\ast} (N) \simeq \varinjlim_{I^{\circ}} \mu_j^{\ast}
m_{j\ast} \mu_{j\ast} (N){\quoi}.\tag{8.5.8.2}
\end{equation*}\etiquette[8.5.8.2]{VI.eq8.5.8.2}
\end{enonce*}

图 \eqhyperref[VI.eq8.5.8.1]{8.5.8.1} 的交换性直接由定义得到
（\hyperref[VI.subsubsec8.2.8]{8.2.8} 及 \hyperref[VI.subsubsec8.5.1]{8.5.1}）。有
$\underleftarrow{m}_{\ast}(N)\simeq Q^{\ast}Q_{\ast}\underleftarrow{m}_{\ast}(N)$，而由
\eqhyperref[VI.eq8.5.8.1]{8.5.8.1} 的交换性，有
$\underleftarrow{m}_{\ast}(N)\simeq Q^{\ast}m^{\sim}_{\ast}Q_{\ast}(N)$。
因此有同构 $\underleftarrow{m}_{\ast}(N)\simeq Q^{\ast}m^{\sim}_{\ast}(j\mapsto\mu_{j\ast}(N))$。
函子 $m_{\ast}^{\sim}$ 就是函子 $\Hom_{I^{\circ}}(I^{\circ},m_{\ast})$
（\hyperref[VI.prop7.4.10]{7.4.10}），因而有典范同构
$\underleftarrow{m}_{\ast}(N)\simeq Q^{\ast}(j\mapsto m_{j\ast}\mu_{j\ast}(N))$。
公式 \eqhyperref[VI.eq8.5.8.2]{8.5.8.2} 于是由 \hyperref[VI.prop8.5.2]{8.5.2} 得到。

\begin{enonce*}{命题 8.5.9}\etiquette[8.5.9]{VI.prop8.5.9}
在 \hyperref[VI.prop8.5.8]{8.5.8} 的假设与记号下，再假设对 $I$ 的每个对象 $i$，函子
$m_{i\ast}:F_i\to G_i$ 都与滤过归纳极限交换，且对 $I$ 的每个态射 $f:i\to j$，函子
$f_{\ast}:F_i\to F_j$ 都与滤过归纳极限交换。则函子
$\underleftarrow{m}_{\ast}:\underleftarrow{F}\to\underleftarrow{G}$ 与滤过归纳极限交换，且对
$\Top(F)$ 的每个对象 $(i\mapsto M_i)$，有典范同构
\begin{equation*}
\underleftarrow{m}_{\ast} Q^{\ast} (i \mapsto M_i) \simeq
\varinjlim_{I^{\circ}} \mu^{\ast}_j m_{j\ast} (M_j){\quoi}.\tag{8.5.9.1}
\end{equation*}\etiquette[8.5.9.1]{VI.eq8.5.9.1}
\end{enonce*}

第一项\pageoriginale{}断言由 \eqhyperref[VI.eq8.5.8.2]{8.5.8.2} 及
\hyperref[VI.cor8.5.6]{8.5.6} 得到。由 \eqhyperref[VI.eq8.5.8.2]{8.5.8.2} 及
\eqhyperref[VI.eq8.5.3.1]{8.5.3.1}，有函子性的同构
$$
\underleftarrow{m}_{\ast} Q^{\ast} (i \mapsto M_i) \simeq
\varinjlim_i \mu^{\ast}_{i} m_{i\ast} (\varinjlim_{f
  : j \to i} f_{\ast} (M_{j})){\quoi}.
$$
利用 $m_{i\ast}$ 与滤过归纳极限的交换性以及同构
$m_{i\ast}f_{\ast}\simeq f_{\ast}m_{j\ast}$（\hyperref[VI.subsubsec7.1.6]{7.1.6}），得到
$$
\underleftarrow{m}_{\ast} Q^{\ast} (i \mapsto M_i) \simeq
\varinjlim_i \mu^{\ast}_{i} (\varinjlim_{f : j \to
  i} f_{\ast} m_{j\ast} (M_j)){\quoi}.
$$
由于函子 $\mu_i^{\ast}$ 与归纳极限交换，得到
$$
\underleftarrow{m}_{\ast}Q^{\ast}(i\mapsto M_{i})\simeq \varinjlim_{i}
\varinjlim_{f:j\to i} \mu^{\ast}_{i}f_{\ast}m_{j\ast}(M_{j}){\quoi}.
$$
最后这个对象可以解释为在范畴 $\text{Fl}(I)^{\circ}$ 上的归纳极限，其中
$\text{Fl}(I)$ 是 $I$ 的态射范畴。设 $\phi:I\to\text{Fl}(I)$ 是将 $I$ 的每个对象 $i$ 映为
$i$ 的恒等态射的函子。函子 $\phi^{\circ}:I^{\circ}\to\text{Fl}(I^{\circ})$ 是共尾的，因而
（I \hyperref[I.sec8.1]{8.1}）有典范同构
$$
\underleftarrow{m}_{\ast} Q^{\ast} (i \mapsto M_i) \simeq
\varinjlim_{I^{\circ}} \mu^{\ast}_{j} m_{j\ast}(M_j){\quoi}.
$$

\begin{enonce*}{推论 8.5.10}\etiquette[8.5.10]{VI.cor8.5.10}
在 \hyperref[VI.prop8.5.9]{8.5.9} 的条件下，对 $I$ 的每个对象 $i$ 及 $F_i$ 的每个对象
$M_i$，有典范同构
\begin{equation*}
\underleftarrow{m}_{\ast} \mu^{\ast}_i (M_i) \simeq
\varinjlim_{f:j\to i} \mu^{\ast}_{j} m_{j\ast}f^{\ast}
(M_{i}){\quoi}.\tag{8.5.10.1}
\end{equation*}\etiquette[8.5.10.1]{VI.eq8.5.10.1}
\end{enonce*}

证明与推论 \hyperref[VI.cor8.5.5]{8.5.5} 的证明类似。

\begin{enonce*}[remark]{注 8.5.11}\etiquette[8.5.11]{VI.rem8.5.11}
回顾凝聚拓扑斯之间的凝聚态射的正像函子与滤过归纳极限交换
（\hyperref[VI.thm5.1]{5.1}）。因此，当所考虑的纤维化拓扑斯为凝聚拓扑斯，且所考虑的纤维化拓扑斯态射为凝聚态射时，命题
\hyperref[VI.prop8.5.3]{8.5.3} 至 \hyperref[VI.cor8.5.7]{8.5.7} 以及
\hyperref[VI.prop8.5.9]{8.5.9}、\hyperref[VI.cor8.5.10]{8.5.10} 均适用。
\end{enonce*}


% END INCLUDED SGA4_English_current_buildable_source_20260906/source/06/drafts/SGA4_VI_LINES_6611_6839_English.texfrag


% BEGIN INCLUDED SGA4_English_current_buildable_source_20260906/source/06/drafts/SGA4_VI_LINES_6840_6915_English.texfrag
\subsection{环化纤维化 $𝒰$-拓扑斯}\etiquette[8.6]{VI.subsec8.6}

\subsubsection{}\etiquette[8.6.1]{VI.subsubsec8.6.1}

若一个纤维化 $𝒰$-拓扑斯配备了总景（\hyperref[VI.def7.4.1]{7.4.1}）上的环层，则称之为\emph{环化}纤维化拓扑斯。
设 $p : F \to I$ 是纤维化 $𝒰$-拓扑斯。选取一个双分裂\pageoriginale{}（\hyperref[VI.subsubsec7.1.4]{7.1.4}）。
由 \hyperref[VI.prop7.4.7]{7.4.7}，给定 $F$ 上的环层等价于对 $I$ 的每个对象 $i$ 给定
$F_i$ 的一个环 $A_i$，并对每个态射 $f : i \to j$ 给定态射
$\phi_f : A_j \to f_{\ast} A_i$，其中族 $\phi_f$ 满足 \loccit{} 中明确写出的相容性条件。
因此，拓扑斯态射 $f. : F_i \to F_j$ 是分别以 $A_i$、$A_j$ 环化的拓扑斯之间的态射（IV 11）；
$F$ 上的环化结构连同态射 $f$ 的选择给出一个从 $I$ 到环化拓扑斯的 2-范畴的伪函子。
反之，给定这样的伪函子，可以重构出产生它的环化纤维化 $𝒰$-拓扑斯。

\subsubsection{}\etiquette[8.6.2]{VI.subsubsec8.6.2}

与 \hyperref[VI.subsec8.1]{8.1} 一样，可以定义环化 $𝒰$-拓扑斯
$(F_i, A_i, i \in I)$ 的投射极限概念。我们假定已经作了这一直接推广，并将
\hyperref[VI.subsec8.1]{8.1} 中引入的语言自由地用于这一情形。若存在，这个投射极限是一个环化
$𝒰$-拓扑斯 $(H, B)$，且对 $I$ 的每个对象 $i$，有环化拓扑斯态射
$\mu_i : (H, B) \to (F_i, A_i)$。若底层的（未环化的）纤维化拓扑斯 $F$ 有投射极限
$\limtop{I} F_i$，且环的归纳系 $i \mapsto \mu_i^{\ast}(A_i)$ 在
$\limtop{I} F_i$ 的环范畴中有归纳极限 $\varinjlim_{I^{\circ}}\mu_i^{\ast} A_i$
（当 $I$ 小时总是如此），则环化拓扑斯
（$\limtop{I} F_i$，$\varinjlim_{I^{\circ}} \mu^{\ast}_i A_i$）显然是所考虑的环化纤维化拓扑斯的投射极限。

\subsubsection{}\etiquette[8.6.3]{VI.subsubsec8.6.3}

\hyperref[VI.subsec8.5]{8.5} 中对集合层建立的公式，对阿贝尔层也成立。只要将其中的逆像函子解释为模层意义下的逆像函子
（IV \hyperref[IV.def1.1]{1.1}），这些公式对模层也成立。验证留给读者。

\subsubsection{}\etiquette[8.6.4]{VI.subsubsec8.6.4}

设 $(p : F \to I, A)$ 是余滤过范畴上的环化纤维化 $𝒰$-拓扑斯。若对每个态射
$f : i \to j$，态射 $f. : (F_i, A_i) \to (F_j, A_j)$ 都是右平坦（\resp{}左平坦）的环化拓扑斯态射
（\hyperrefext[V][V.def1.8]{1.8}），则称 $(p : F \to I, A)$ 为右平坦（\resp{}左平坦）的环化纤维化
$𝒰$-拓扑斯。假设 $I$ 本质小。则环化拓扑斯态射
$Q : (\limtop{I} F_i,
\varinjlim_{I^{\circ}} \mu^{\ast}_i (A_i)) \to
(\Hom_{I^{\circ}}(I^{\circ}, F'), A)$\pageoriginale{}
（\hyperref[VI.subsubsec8.5.1]{8.5.1}）为右平坦（\resp{}左平坦），这是定义及
\hyperref[VI.prop8.5.2]{8.5.2} 的推论；此结论不要求过渡态射 $f.$ 满足任何平坦性假设。
若进一步 $(p : F \to I, A)$ 是右平坦（\resp{}左平坦）的环化纤维化 $𝒰$-拓扑斯，则态射
$\mu_i : (\limtop{I} F_i,
\varinjlim_{I^{\circ}} \mu^{\ast}_i A_i) \to (F_i, A_i)$ 为右平坦（\resp{}左平坦）。
这由平坦模的逆像仍然平坦（\hyperrefext[V][V.cor1.7.1]{1.7.1}）以及平坦模的滤过归纳极限仍然平坦这两个事实得到。
同样的论证证明下列事实：若 $m : (F_i, A_i, i \in I) \to (G_i, A_i, i \in I)$
是本质小余滤过范畴上的环化纤维化 $𝒰$-拓扑斯之间的笛卡儿态射，并且对每个对象
$i \in \ob (I)$，态射 $m_i$ 为右平坦（\resp{}左平坦），则由 $m$ 取投射极限而导出的态射
$\underleftarrow{m}$ 为右平坦（\resp{}左平坦）。

% END INCLUDED SGA4_English_current_buildable_source_20260906/source/06/drafts/SGA4_VI_LINES_6840_6915_English.texfrag


% BEGIN INCLUDED SGA4_English_current_buildable_source_20260906/source/06/drafts/SGA4_VI_LINES_6916_7181_English.texfrag
\subsection{拓扑斯投射极限上层的上同调}\etiquette[8.7]{VI.subsec8.7}

\subsubsection{}\etiquette[8.7.1]{VI.subsubsec8.7.1}

本节中，$I$ 表示小余滤过范畴，$(p:F\to I,A)$ 及 $(q:G\to I,B)$ 表示两个环化纤维化
$𝒰$-拓扑斯，$m:(p:F\to I,A)\to(q:G\to I,B)$ 表示环化纤维化 $𝒰$-拓扑斯态射。
假设对 $I$ 的每个态射 $f:i\to j$，模函子
$f_{\ast}:\Mod(F_i,A_i)\to\Mod(F_j,A_j)$ 的导出函子 $R^nf_{\ast}$，$n\in\NN$，都与滤过归纳极限交换，
且对每个对象 $i\in I$，模函子\footnote{为明确起见，我们取左模。}
$m_{i\ast}:(F_i,A_i)\to(G_i,B_i)$ 的导出函子 $R^nm_{i\ast}$ 都与滤过归纳极限交换。
回顾当范畴 $F_i,G_i$，$(i\in\ob(I))$，为\emph{凝聚}，且态射
$f:F_i\to F_j$，$(f:i\to j\in \text{Fl}(I))$，及
$m_i:F_i\to G_i$，$(i\in I)$，为\emph{凝聚}时，这些假设均成立（\hyperref[VI.cor5.2]{5.2}）。
采用 \hyperref[VI.subsubsec8.5.1]{8.5.1} 的记号。此外，以 $\underrightarrow{A}$
（\resp $\underrightarrow{B}$）表示层 $\varinjlim\mu^{\ast}_i(A_i)$
（\resp $\varinjlim\mu^{\ast}_i(B_i)$）。最后，对模，我们用记号 $\mu^{-1}_i$ 表示阿贝尔层意义下的逆像，
而将 $\mu_i^{\ast}$ 留作模意义下的逆像（IV 11）。

\begin{enonce*}{引理 8.7.2}\etiquette[8.7.2]{VI.lem8.7.2}
设\pageoriginale{} $j\mapsto M_j$ 是 $\Top(F)$ 的内射 $A$-模。则
$Q^{\ast}(j\mapsto M_j)$ 是对 $\underleftarrow{m}_{\ast}$ 无上同调的 $A$-模，且对每个
$i\in\ob(I)$，$M_i$ 是松弛的。
\end{enonce*}

设 $i\in\ob(I)$，并设 $e_i$ 是 $F_i$ 的终对象。纤维化 $𝒰$-拓扑斯
$F_{/e_i}\to I_{/i}$ 由 $F\to I$ 经基变换 $I_{/i}\to I$ 得到。
由于层 $j\mapsto M_j$ 是内射的，它是松弛的（\hyperrefext[V][V.subsec4.6]{4.6}）。
它在局部化拓扑斯 $F_{/e_i}$ 上的限制是松弛的（\hyperrefext[V][V.subsec4.12]{4.12}）。
从 $F_{/e_i}$ 到 $F_i$ 的限制函子是某个拓扑斯态射的正像函子
（\hyperref[VI.lem7.4.12]{7.4.12}），因而将松弛模映为松弛模（\hyperrefext[V][V.prop5.2]{5.2}）。
所以 $M_i$ 是松弛的。

证明引理的第一项断言。令 $N'=Q^{\ast}(j\mapsto M_j)$。
由 \eqhyperref[VI.eq8.5.3.1]{8.5.3.1}，对每个 $i\in\ob(I)$ 有
$$
\mu_{i\ast}(N') = \mu_{i\ast} Q^{\ast} (j \mapsto M_j) \simeq
\varinjlim_{f : j \to i} f_{\ast} (M_j){\quoi}.
$$
利用 \hyperref[VI.subsubsec8.7.1]{8.7.1} 的假设，可见层 $N'$ 具有下列性质：

(P)~对每个对象 $i\in\ob(I)$，$\mu_{i\ast}(N')$ 对 $m_{i\ast}$ 无上同调，且对
$I$ 的每个态射 $f:i\to j$，$\mu_{i\ast}(N')$ 对 $f_{\ast}$ 无上同调。

只需证明每个具有性质 (P) 的对象 $N'$ 都对 $\underleftarrow{m}_{\ast}$ 无上同调。
由于内射对象具有性质 (P)，只需验证下列两个性质（\hyperrefext[V][V.prop0.4]{0.4}）：
对每个正合列 $0\to N'\to N\to N''\to0$，其中 $N'$ 及 $N$ 具有性质 (P)，
(a) $N''$ 具有性质 (P)，且 (b) $m_{\ast}(N)\to m_{\ast}(N'')$ 是满态射。
(a) 的验证是直接的。验证 (b)。令 $K=\coker(Q_{\ast}(N')\to Q_{\ast}(N))$，
于是 $K$ 是截面 $(i\mapsto K_i=\coker(\mu_{i\ast}(N')\to\mu_{i\ast}(N)))$。
对每个 $f:i\to j$，有交换图
\[
\xymatrix{
0  \ar[r] & f_{\ast}\mu_{i\ast}(N') \ar[r] & f_{\ast}\mu_{i\ast}(N)\ar[r]
 & f_{\ast}(K_i) \ar[r]& 0\\
0  \ar[r] & \mu_{j\ast}(N')\ar[u] \ar[r] & \mu_{j\ast}(N) \ar[u] \ar[r] & K_j
\ar[r] \ar[u]&0.
}
\]
由于\pageoriginale{} $\mu_{i\ast}(N')$ 对 $f_{\ast}$ 无上同调，上行正合。
由于 $i\mapsto\mu_{i\ast}(N')$ 及 $i\mapsto\mu_{i\ast}(N)$ 是笛卡儿截面，前两个竖直态射是同构。
因此，典范态射 $K_j\to f_{\ast}(K_i)$ 是同构，$i\mapsto K_i$ 是笛卡儿截面。
所以 $K$ 具有 $Q_{\ast}(K')$ 的形式；由于 $Q_{\ast}$ 全忠实，典范态射
$Q_{\ast}(N)\to K$ 及 $K\to Q_{\ast}(N'')$ 分别具有 $Q_{\ast}(u)$ 及 $Q_{\ast}(v)$ 的形式。
由于 $Q^{\ast}$ 正合，且 $Q^{\ast}Q_{\ast}$ 同构于恒等函子，$v$ 是同构，故 $Q_{\ast}(v)$ 是同构。
所以列 $0\to Q_{\ast}(N')\to Q_{\ast}(N)\to Q_{\ast}(N'')\to0$ 正合，
因而对 $I$ 的每个对象 $i$，列 $0\to\mu_{i\ast}(N')\to\mu_{i\ast}(N)\to\mu_{i\ast}(N'')\to0$ 正合。
由于 $\mu_{i\ast}(N')$ 对 $m_{i\ast}$ 无上同调，列
$0\to m_{i\ast}\mu_{i\ast}(N')\to m_{i\ast}\mu_{i\ast}(N)\to m_{i\ast}\mu_{i\ast}(N'')\to0$ 正合。
此外，有同构
$\underleftarrow{m}_{\ast}\simeq \varinjlim_{I^{\circ}}\mu^{-1}_jm_{j\ast}\mu_{j\ast}$
\eqhyperref[VI.eq8.5.8.2]{8.5.8.2}。函子 $\mu^{-1}_j$ 正合，且滤过归纳极限正合。
因此，列 $0\to\underleftarrow{m}_{\ast}(N')\to\underleftarrow{m}_{\ast}(N)\to\underleftarrow{m}_{\ast}(N'')\to0$ 正合。

\begin{enonce*}{定理 8.7.3}\etiquette[8.7.3]{VI.thm8.7.3}
记号与假设同 \hyperref[VI.subsubsec8.7.1]{8.7.1}。对每个整数 $n\in\NN$，导出函子
$R^n\underleftarrow{m}_{\ast}$ 与滤过归纳极限交换，并且对每个 $A$-模
$j\mapsto M_j$，有函子性的同构
\begin{equation*}
R^{n} \underleftarrow{m}_{\ast} Q^{\ast} (j \mapsto M_j) \simeq
\varinjlim_{I^{\circ}}
\mu^{\ast}_j R^n m_{j\ast} (M_j)\tag{8.7.3.1}{\quoi}.
\end{equation*}\etiquette[8.7.3.1]{VI.eq8.7.3.1}
\end{enonce*}

有 $Q^{-1}(A)=\underrightarrow{A}$，故模的逆像同构于阿贝尔层的逆像。
由 \hyperref[VI.eq8.5.8.2]{8.5.8.2}，对每个截面 $j\mapsto M_j$，有典范同构
$\varinjlim_{I^{\circ}}\mu^{-1}_j(M_j)\simeq \varinjlim_{I^{\circ}}\mu^{\ast}_j(M_j)$。
此外，在 \hyperref[VI.subsubsec8.7.1]{8.7.1} 的假设下，有典范同构
\eqhyperref[VI.eq8.5.9.1]{8.5.9.1}
$\underleftarrow{m}_{\ast}Q^{\ast}(j\mapsto M_j)\simeq \varinjlim_{I^{\circ}}\mu^{-1}_jm_{j\ast}(M_j)$。
由于对每个内射的 $j\mapsto M_j$，$M_j$ 都是松弛的，且 $Q^{\ast}(j\mapsto M_j)$
对 $\underleftarrow{m}_{\ast}$ 无上同调（\hyperref[VI.lem8.7.2]{8.7.2}），有同构
$R^n\underleftarrow{m}_{\ast}(Q^{\ast}(j\mapsto M_j))\simeq \varinjlim_{I^{\circ}}\mu^{-1}_jR^nm_{j\ast}(M_j)$，
结合前述结论即得公式 \eqhyperref[VI.eq8.7.3.1]{8.7.3.1}。
第一项断言直接由公式 \eqhyperref[VI.eq8.7.3.1]{8.7.3.1} 得到。

\begin{enonce*}{推论 8.7.4}\etiquette[8.7.4]{VI.cor8.7.4}
对\pageoriginale{} $n\in\ZZ$，有函子性的同构
\begin{equation*}
R^n \underleftarrow{m}_{\ast} (N) \simeq \varinjlim_{I^{\circ}}
\mu^{\ast}_j R^n m_{j\ast} \mu_{j\ast} (N){\quoi}.\tag{8.7.4.1}
\end{equation*}\etiquette[8.7.4.1]{VI.eq8.7.4.1}
\end{enonce*}

事实上，$N$ 同构于 $Q^{\ast}(j\mapsto\mu_{j\ast}(N))$。

\begin{enonce*}{推论 8.7.5}\etiquette[8.7.5]{VI.cor8.7.5}
对 $I$ 的每个对象 $i$、每个整数 $n$ 及每个 $A_i$-模 $M_i$，有函子性的同构
\begin{equation*}
R^n \underleftarrow{m}_{\ast} \mu^{\ast}_i (M_i) \simeq
\varinjlim_{f : j \to i} \mu^{\ast}_j R^n m_{j\ast} f^{\ast}
(M_i){\quoi}.\tag{8.7.5.1}
\end{equation*}\etiquette[8.7.5.1]{VI.eq8.7.5.1}
\end{enonce*}

作基变换 $I_{/i}\to I$ 后，可以假设 $i$ 是 $I$ 的终对象
（\hyperref[VI.lem8.2.1]{8.2.1}）。此时对象 $\mu_i^{\ast}(M_i)$ 同构于
$Q^{\ast}((f:j\to i)\mapsto f^{\ast}(M_i))$（\cf \hyperref[VI.cor8.5.5]{8.5.5}）。

\begin{enonce*}{推论 8.7.6}\etiquette[8.7.6]{VI.cor8.7.6}
设 $i$ 是 $I$ 的对象，$n$ 是整数。有关于截面 $j\mapsto M_j$ 函子性的同构：
\begin{equation*}
R^n \mu_{i\ast} Q^{\ast} (j \mapsto M_j) \simeq \varinjlim_{f :
  j \to i} R^n f_{\ast} (M_j){\quoi}.\tag{8.7.6.1}
\end{equation*}\etiquette[8.7.6.1]{VI.eq8.7.6.1}
有关于 $\underrightarrow{A}$-模 $N$ 函子性的同构：
\begin{equation*}
R^n \mu_{i\ast} (N) \simeq \varinjlim_{f:j\to i}
R^n f_{\ast} \mu_{j\ast} (N){\quoi}.\tag{8.7.6.2}
\end{equation*}\etiquette[8.7.6.2]{VI.eq8.7.6.2}
有关于 $A_i$-模 $M_i$ 函子性的同构：
\begin{equation*}
R^n \mu_{i\ast} \mu^{\ast}_i (M_i) \simeq \varinjlim_{f:j\to i}
R^n f_{\ast} f^{\ast} (M_i){\quoi}.\tag{8.7.6.3}
\end{equation*}\etiquette[8.7.6.3]{VI.eq8.7.6.3}
函子 $R^n\mu_{i\ast}$ 与滤过归纳极限交换。
\end{enonce*}

作基变换 $I_{/i}\to I$，即归结到 $i$ 是 $I$ 的终对象的情形。此时，公式
\eqhyperref[VI.eq8.7.6.1]{8.7.6.1} 及 \eqhyperref[VI.eq8.7.6.2]{8.7.6.2} 分别是公式
\eqhyperref[VI.eq8.7.3.1]{8.7.3.1}、\eqhyperref[VI.eq8.7.4.1]{8.7.4.1} 及
\eqhyperref[VI.eq8.7.5.1]{8.7.5.1} 的特例：取 $G$ 为纤维 $F_i$ 的常纤维化拓扑斯，
取 $m$ 为态射 $(f:F_j\to F_i, f\in\ob(I_{/i}))$ 即可得到。
最后一项断言由 \eqhyperref[VI.eq8.7.6.2]{8.7.6.2} 并结合 \hyperref[VI.cor8.5.4]{8.5.4} 得到。

\begin{enonce*}{推论 8.7.7}\etiquette[8.7.7]{VI.cor8.7.7}
设\pageoriginale{} $i$ 是 $I$ 的对象，$X$ 是 $F_i$ 的对象，使得函子
$H^n(X,-)$，$n\in\NN$，与滤过归纳极限交换。对每个 $n$，有关于 $A_i$-模 $M_i$ 函子性的典范同构：
\begin{equation*}
H^n (\mu_{i}^{\ast}(X), \mu^{\ast}_i (M_i)) \simeq \varinjlim_{f:j\to i}
H^n (f^{\ast} (X), f^{\ast} (M_i))\tag{8.7.7.1}
\end{equation*}\etiquette[8.7.7.1]{VI.eq8.7.7.1}
且函子 $H^n(\mu_i^{\ast}(X),-)$ 与滤过归纳极限交换。特别地，若函子
$H^n(F_i,-)$ 与滤过归纳极限交换，则函子 $H^n(\underleftarrow{F},-)$ 与滤过归纳极限交换，
且有关于 $A_i$-模 $M_i$ 函子性的典范同构：
\begin{equation*}
H^n (\underleftarrow{F}, \mu_{i}^{\ast} (M_i)) \simeq \varinjlim_{f:j\to i}
H^n (F_j, f^{\ast} (M_i)){\quoi}.\tag{8.7.7.2}
\end{equation*}\etiquette[8.7.7.2]{VI.eq8.7.7.2}
\end{enonce*}

\noindent
由 \hyperrefext[V][V.prop5.3]{5.3}，有两个谱序列
\begin{align*}
'E_{2}^{p, q} &= \varinjlim_{f:j\to i} H^p(X, R^q f_{\ast}
f^{\ast} (M_i)) \Rightarrow \varinjlim_{f:j\to i}
H^{p+q}(f^{\ast} (X), f^{\ast} (M_i))\\
''E_{2}^{p, q} &=  H^p(X, R^q \mu_{i\ast} \mu^{\ast}_i (M_i))
\Rightarrow H^{p+q} (\mu^{\ast}_i (X), \mu^{\ast}_i (M_i)){\quoi},
\end{align*}
且由 \eqhyperref[VI.eq8.5.7.1]{8.5.7.1}，这两个谱序列之间有一个态射。
由于 $H^{p+q}(X,-)$ 与滤过归纳极限交换，由 \eqhyperref[VI.eq8.7.6.2]{8.7.6.2}，
这个谱序列态射在 $E_2^{p,q}$ 项上是同构。因此它在收敛对象上诱导同构，给出
\eqhyperref[VI.eq8.7.7.1]{8.7.7.1}。对每个 $\underrightarrow{A}$-模 $N$，有谱序列
（\hyperrefext[V][V.prop5.3]{5.3}）
$$
E^{p, q}_2 = H^{p} (X, R^q \mu_{i\ast} (N)) \Rightarrow H^{p+q}
(\mu^{\ast}_i (X), N){\quoi}.
$$
于是，函子 $H^n(\mu_i^{\ast}(X),-)$ 与滤过归纳极限交换的断言，由函子
$H^n(X,-)$ 及 $R^n\mu_{i\ast}$ 的类似性质得到（\hyperref[VI.cor8.7.6]{8.7.6}）。

\begin{enonce*}[remark]{注 8.7.8}\etiquette[8.7.8]{VI.rem8.7.8}
当拓扑斯 $F_j$，$j\in\ob(I)$，为凝聚，且过渡态射为凝聚时，若 $X$（\resp $F_i$）
为凝聚，则 \hyperref[VI.cor8.7.7]{8.7.7} 中对 $X$（\resp $F_i$）所作的假设成立
（\hyperref[VI.cor5.2]{5.2}）。此外，我们知道在此情形下 $\mu_i^{\ast}(X)$
（\resp $\underleftarrow{F}$）为凝聚（\hyperref[VI.thm8.3.13]{8.3.13}）。
\end{enonce*}

\begin{enonce*}{推论 8.7.9}\etiquette[8.7.9]{VI.cor8.7.9}
设\pageoriginale{} $i$ 是 $I$ 的对象，$M_i$ 是左 $A_i$-模，$L_i$ 是右 $A_i$-模，并有形如
$$
P_{i, k} \to P_{i, k-1} \to \cdots \to P_{i, 0} \to L_i \to 0
$$
的消解，其中对每个整数 $k$，$P_{i,k}$ 同构于形如 $A_{i|X}$ 的对象的有限直和
\emph{（IV \hyperref[IV.prop11.3.3]{11.3.3}）}，且 $X$ 满足
\hyperref[VI.cor8.7.7]{8.7.7} 的假设。若除 \hyperref[VI.subsubsec8.7.1]{8.7.1} 的假设外，
纤维化拓扑斯 $F$ 还是右平坦的（\hyperref[VI.subsec8.6]{8.6}），则对 $n\leq k-1$，有典范同构
\begin{equation*}
\Ext^n_{\underrightarrow{A}}(\underleftarrow{F}, \mu^{\ast}_i(L_i),
\mu^{\ast}_i (M_i)) \simeq \varinjlim_{f : j \to i}
\Ext^{n}_{A_j} (F_j, f^{\ast}(L_i), f^{\ast} (M_i)){\quoi}.\tag{8.7.9.1}
\end{equation*}\etiquette[8.7.9.1]{VI.eq8.7.9.1}
\end{enonce*}

\noindent
以 $P_{i,.}$ 表示 $L_i$ 的消解。由于 $F$ 右平坦，对每个 $f:j\to i$，
$f^{\ast}(P_{i,.})$ 是 $f^{\ast}(L_i)$ 的消解，且 $\mu_i^{\ast}(P_{i,.})$ 是
$\mu_i^{\ast}(L_i)$ 的消解（\hyperref[VI.subsec8.6]{8.6}）。这些消解给出两个分别收敛于
\eqhyperref[VI.eq8.7.9.1]{8.7.9.1} 两端的谱序列，以及两个谱序列之间的一个态射。
由 \eqhyperref[VI.eq8.7.7.1]{8.7.7.1}，此态射在 $E_1^{p,q}$ 项上是同构。因此它在收敛对象上诱导同构。


% END INCLUDED SGA4_English_current_buildable_source_20260906/source/06/drafts/SGA4_VI_LINES_6916_7181_English.texfrag


% BEGIN INCLUDED SGA4_English_current_buildable_source_20260906/source/06/drafts/SGA4_VI_LINES_7182_7367_English.texfrag
\section{附录：点存在的判据}\etiquette[9]{VI.sec9}

\centerline{P. \textsc{Deligne} \textit{著}}

\medskip
\begin{enonce*}{命题 9.0}\etiquette[9.0]{VI.prop9.0}
每个局部凝聚拓扑斯 $S$ 都有足够多的点。
\end{enonce*}

由于问题在 $S$ 上是局部的，可以假设拓扑斯 $S$ 由景 $\Ss$ 定义，其中有限投射极限可表，
且每个覆盖 $f_i:U_i\to U$ 都有有限子覆盖。只需证明：若 $f:\Ff\to\Gg$ 不是单态射，
则存在 $\Ss$ 的一个点 $x$，使 $f_x$ 不是单射。由假设，存在 $U\in\ob\Ss$ 及
$\ss,\ss'\in\Ff(U)$，使得 $\ss\ne\ss'$ 且 $f(\ss)=f(\ss')$。
将 $\Ss$ 换为 $\Ss/U$，即归结到下列引理。

\begin{enonce*}{引理 9.1}\etiquette[9.1]{VI.lem9.1}
若\pageoriginale{} $\ss$ 与 $\ss'$ 是满足上述假设的景 $\Ss$ 上某个层 $\mathfrak F$ 的两个不同整体截面，
则存在 $\Ss$ 的点 $x$，使得 $\ss_x\ne\ss'_x$。
\end{enonce*}

设 $P=(U_i)_{i\in I}$ 是 $\Ss$ 中以滤过有序集 $I$ 为指标的投射系。对 $\Ss$ 上的每个层
$\Hh$，令 $P(\Hh)=\varinjlim\Hh(U_i)$；对每个 $V\in\Ob\Ss$，令
$P(V)=\varinjlim\Hom(U_i,V)$。函子 $P(\Hh)$ 及 $P(V)$ 与纤维积交换。
若函子 $P(V)$ 将覆盖映为满射映射族，则它是从点拓扑斯到 $\Ss$ 的景态射，而
$P(\Hh)$ 是相应的纤维函子。

若 $P=(U_i)_{i\in I}$ 及 $Q=(V_j)_{j\in J}$ 是 $\Ss$ 中的两个投射系，
当 $I$ 是 $J$ 的子集并带有诱导的序，且 $P$ 是 $Q$ 在 $I$ 上的限制时，称 $Q$ 加细 $P$。
此时有从 $P(\Hh)$ 到 $Q(\Hh)$ 及从 $P(V)$ 到 $Q(V)$ 的函子态射。

对每个如上的 $P$，以 $\ss_P$ 及 $\ss'_P$ 表示 $\ss$ 及 $\ss'$ 在 $P(\Ff)$ 中的像。
在下列引理中，取 $P$ 为以 $\{0\}$ 为指标且由 $\Ss$ 的终对象构成的投射系，
即可得到引理 \hyperref[VI.lem9.1]{9.1}。

\begin{enonce*}{引理 9.2}\etiquette[9.2]{VI.lem9.2}
对每个满足 $\ss_P\ne\ss'_P$ 的如上 $P$，存在 $P$ 的加细 $Q$，满足
$\ss_Q\ne\ss'_Q$，且对每个覆盖 $f_i:V_i\to V$ 及每个 $f\in Q(V)$，
$f$ 属于某个 $Q(V_i)$ 的像。
\end{enonce*}

先证明：

\begin{enonce*}{引理 9.3}\etiquette[9.3]{VI.lem9.3}
设 $P=(U_i)_{i\in I}$ 满足引理 \hyperref[VI.lem9.2]{9.2} 的假设，
$f_i:V_i\to V$ 是 $\Ss$ 中的有限覆盖，$f\in P(V)$。
存在 $P$ 的加细 $Q$，满足 $\ss_Q\ne\ss'_Q$，且 $f$ 在 $Q(V)$ 中的像属于某个
$Q(V_i)$ 的像。
\end{enonce*}

存在 $i_{\circ}\in I$ 及 $f'\in\Hom(U_{i_{\circ}},V)$，它们定义 $f$。
对 $i\geq i_{\circ}$，令 $U_{i,k}=\fprod{U_i}{V_k}{V}$，此纤维积由 $f'$ 定义。
由于 $(V_k)$ 是 $V$ 的覆盖，$(U_{i,k})$ 是 $U_i$ 的覆盖；下列箭头为单射（对 $i>i_{\circ}$）：
$$
\Ff(U_i)\to\prod_k\Ff(U_{i,k}).
$$
取\pageoriginale{}极限，并利用\emph{有限}积与滤过归纳极限的交换性，可见
$$
P(\Ff)\to\prod_k\varinjlim\Ff(U_{i,k})
$$
是单射。因此，存在 $k$，使得 $\ss_P$ 与 $\ss'_P$ 在
$\varinjlim_{i>i_{\circ}}\Ff(U_{i,k})$ 中的像不同。令
$I_{\circ}=\{i\mid i\geq i_{\circ}\text{ 且 }i\in I\}$ 及
$J=I\coprod I_{\circ}$。在 $J$ 上规定次序：$j'\leq j''$ 当且仅当
$j'$ 及 $j''$ 在 $I$ 中的像满足 $j'\leq j''$，且不出现 $j'\in I_{\circ}$ 而 $j''\in I$ 的情形。
以 $J$ 为指标的 $U_i$ 及 $U_{i,k}$ 构成 $P$ 的一个加细 $Q$，使得
$\ss_Q\ne\ss'_Q$，且 $f$ 在 $Q(V)$ 中的像属于 $Q(V_k)$ 的像。

\begin{enonce*}{引理 9.4}\etiquette[9.4]{VI.lem9.4}
设 $P=(U_i)_{i\in I}$ 满足引理 \hyperref[VI.lem9.2]{9.2} 的假设。
存在 $P$ 的加细 $Q$，使得 $\ss_Q\ne\ss'_Q$，且对 $\Ss$ 的对象 $V$ 的每个有限覆盖
$(V_k)$ 及每个 $f\in P(V)$，$f$ 在 $Q(V)$ 中的像属于某个 $Q(V_k)$ 的像。
\end{enonce*}

设 $E$ 是由 $\Ss$ 的对象 $V$、$V$ 的有限覆盖 $(V_k)$ 及元素 $f\in P(V)$ 组成的三元组的集合。
设 $\leq$ 是 $E$ 的一个良序，$\bar E$ 是向 $E$ 添加一个记为 $\infty$ 的最大元所得的集合。
通过对 $e\in\bar E$ 作超限归纳，我们将定义 $P$ 的加细 $Q_e$，满足
$\ss_{Q_e}\ne\ss'_{Q_e}$，且
\begin{enumerate}
\renewcommand{\theenumi}{\roman{enumi}}
\renewcommand{\labelenumi}{(\theenumi)}
\item 若 $e'\leq e$，则 $Q_e$ 加细 $Q_{e'}$；

\item 若 $e=(V,(V_k),f)$，则 $f$ 在 $Q_e(V)$ 中的像属于某个 $Q_e(V_k)$ 的像。
\end{enumerate}

假设对 $e'<e$ 已经定义了 $Q_{e'}$。若 $e$ 是 $E$ 的第一个元素，令 $Q'_e=P$。
若 $e$ 有前驱 $e-1$，令 $Q'_e=Q_{e-1}$。否则，设 $Q'_e$ 是以
$\bigcup_{e'<e}I_{e'}$ 为指标集、加细所有 $e'<e$ 对应的 $Q_{e'}$ 的投射系。
在此情形下，
$$
Q'_e(\Ff)=\varinjlim_{e'<e}Q_{e'}(\Ff),
$$
所以在每种情形下，\pageoriginale{} $Q'_e$ 都加细 $e'<e$ 对应的 $Q_{e'}$，且满足
$\ss_{Q'_e}\ne\ss'_{Q'_e}$。

将引理 \hyperref[VI.lem9.3]{9.3} 应用于 $Q'_e$ 及 $e=(V,(V_k),f)$，即可得到所需的投射系
$Q_e$（\resp{}若 $e=\infty$，则取 $Q_e=Q'_e$）。

投射系 $Q_{\infty}$ 满足引理 \hyperref[VI.lem9.4]{9.4}。

利用引理 \hyperref[VI.lem9.4]{9.4}，可对 $n$ 归纳地定义 $\Ss$ 中的投射系序列
$Q_n=(U_i)_{i\in I_n}$，使得
\begin{enumerate}
\renewcommand{\theenumi}{\roman{enumi}}
\renewcommand{\labelenumi}{(\theenumi)}
\item $Q_0=P$；

\item $Q_{n+1}$ 加细 $Q_n$；

\item $\ss_{Q_n}\ne\ss'_{Q_n}$；

\item 对每个覆盖 $f_k:V_k\to V$ 及每个 $f\in Q_n(V)$，
$f$ 在 $Q_{n+1}(V)$ 中的像属于某个 $Q_{n+1}(V_k)$ 的像。
\end{enumerate}

以各 $I_n$ 的并为指标集并延拓所有 $Q_n$ 的投射系 $Q$，于是满足引理
\hyperref[VI.lem9.2]{9.2}。证明还表明：

\begin{enonce*}{推论 9.5}\etiquette[9.5]{VI.cor9.5}
设 $\Ss$ 是纤维积可表的景，且 $\Ss$ 中每个覆盖都有有限子覆盖。令
$c=\sup(\aleph_0,\card(\text{Fl }\Ss))$。则存在 $\Ss$ 的一个很稠密的点集 $X$，使得
\begin{enumerate}
\renewcommand{\theenumi}{\roman{enumi}}
\renewcommand{\labelenumi}{(\theenumi)}
\item $\card(X)\leq2^c$；

\item 若 $x\in X$ 且 $U\in\Ob\Ss$，则 $\card(U_x)\leq c$。
\end{enumerate}
\end{enonce*}

\begin{thebibliography}{99}\pageoriginale
\bibitem{sga4.exp.VI.bib.1} P.\textsc{Gabriel} 与 M.\textsc{Zisman}: Homotopy Theory
  and Calculus of Fractions, Ergebnisse der Mathematik, Bd 35.

\bibitem{sga4.exp.VI.bib.2} J. \textsc{Giraud}: M\'ethode de la descente. M\'emoire de la S.M.F.

\bibitem{sga4.exp.VI.bib.3} A. \textsc{Grothendieck} 与 J. \textsc{Dieudonn\'e}: \'El\'ements de
  G\'eom\'etrie Alg\'ebrique I，第 2 版。

\bibitem{sga4.exp.VI.bib.4} A. \textsc{Grothendieck} 与 J. \textsc{Dieudonn\'e}: \'El\'ements de
  G\'eom\'etrie Alg\'ebrique IV I, I.H.E.S.~${\rm n^{\circ}}$

\bibitem{sga4.exp.VI.bib.5} \SGA. 1 VI，A. \textsc{Grothendieck} 著，I.H.E.S.

\bibitem{sga4.exp.VI.bib.6} \SGA. 3 IV 6.3，M. \textsc{Demazure} 著，Lecture Notes
  ${\rm n^{\circ}}$, Springer Verlag.

\bibitem{sga4.exp.VI.bib.7} A. \textsc{Grothendieck}: Sur quelques points d'alg\`ebre
  homologique, Tohoku Math. J.
\end{thebibliography}

\ifx\danslelivre\undefined
\end{document}
\fi

% END INCLUDED SGA4_English_current_buildable_source_20260906/source/06/drafts/SGA4_VI_LINES_7182_7367_English.texfrag


